	\documentclass[14pt]{article}
\usepackage{extsizes}
	\usepackage[frenchb]{babel}
	\usepackage[utf8]{inputenc}  
	\usepackage[T1]{fontenc}
	\usepackage{amssymb}
	\usepackage[mathscr]{euscript}
	\usepackage{stmaryrd}
	\usepackage{amsmath}
	\usepackage{tikz}
	\usepackage[all,cmtip]{xy}
	\usepackage{amsthm}
	\usepackage{varioref}
	\usepackage[ margin=1in]{geometry}
	\geometry{
 a4paper,
 total={170mm,277mm},
 left=20mm,
 top=20mm,
 }
	\usepackage{lmodern}
	\usepackage{hyperref}
	\usepackage{array}
	\usepackage{float}
	\usepackage{easytable}
	 \usepackage{fancyhdr}\usepackage{longtable}
	 \usetikzlibrary{shapes.misc}
	 \newcommand\ang[1]{$#1{}^o$}
\newlength{\taillecellule}
\setlength{\taillecellule}{2cm}
\newcolumntype{C}{@{}>{\centering\arraybackslash}p{\taillecellule}@{}}

	\theoremstyle{remark}
	\newtheorem*{sol}{Solution}
\renewcommand{\theenumi}{\alph{enumi})}
\usepackage{pstricks,multido}
\usepackage{arrayjob}
\usepackage{calc,xlop}
\tikzset{cross/.style={cross out, draw=black, minimum size=2*(#1-\pgflinewidth), inner sep=0pt, outer sep=0pt},
%default radius will be 1pt. 
cross/.default={1pt}}

	\newcounter{n}
	\setcounter{n}{1} 
	\newcommand{\cm}{\textbf{ cm}}
	\newcommand\ds\displaystyle
	\renewcommand{\it}{\item[$ \mathbf{\bullet\ \then}.$]\stepcounter{n} }
	\newenvironment{calcul}{\item[$\mathbf{\then}.$] \stepcounter{n}\hfil $\displaystyle  }{.$\hfil   }
	\pagestyle{fancy}
	\theoremstyle{plain}
	\fancyfoot[C]{\empty} 
	\fancyhead[L]{Bilan}
	\fancyhead[R]{Année de 3e}
	
	
	
	\title{Bilan année de 3e}
	\date{}
	\begin{document}

\section*{Avant-propos}
Voici la correction de la feuille d'exercices  disponible ici : \url{https://leturcqd.github.io/ressources/3e/R\%C3\%A9visions3e/R\%C3\%A9visions3e.pdf}
Elle est disponible à tout instant au lien suivant, qui sera mis à jour : \url{https://leturcqd.github.io/ressources/3e/R\%C3\%A9visions3e/R\%C3\%A9visions3ec.pdf}


Corrigés rédigés dans cette version : \begin{itemize}
\item Partie \ref{calc} : complet.
\item Partie \ref{pbs} de a). à e). et q). à w).
\end{itemize} 
\section{Révision de calcul}\label{calc} 
\begin{itemize}
\it \begin{eqnarray*}
\left(-\frac14\right)+\left(-\frac32\right)+\left(+\frac52\right) & = & \frac{-1}4 + \frac{-3\times 2}{2\times2} + \frac{5\times2}{2\times2} \\
&=& \frac{-1}4 + \frac{-6}4 + \frac{10}4\\
&=& \frac{-1-6+10}4\\
&=& \frac{3}4
\end{eqnarray*} 
ou, plus malin : \begin{eqnarray*}
\left(-\frac14\right)+\left(-\frac32\right)+\left(+\frac52\right) & = & -\frac14 + \frac{-3+5}2\\
& =& - \frac14 + \frac{2}2\\
&=& -\frac14+1\\
&=& -\frac14+\frac44 \\
&=& \frac34
\end{eqnarray*} 
\it \begin{eqnarray*}
\left(-\frac56\right)+\left(+\frac49\right)+\left(-\frac73\right) &= & \frac{-5\times3}{6\times3} + \frac{4\times2}{9\times2} + \frac{-7\times 6}{3\times 6} \\ 
&=& \frac{-15}{18} + \frac{8}{18} +\frac{-42}{18} \\
&=& \frac{-15+8-42}{18}\\
&=& \frac{-49}{18}
\end{eqnarray*}
\it \begin{eqnarray*}(-0,75) + (+4) + \left(-\frac45\right) & =& -0,75 + 4 - 0,8 \\ 
&=& 3,25 - 0,8\\
&=& 2,45
\end{eqnarray*}
\it \begin{eqnarray*}
\left(+\frac47\right)+\left(-\frac{13}{21}\right)+\left(+\frac56\right) &=& \frac{4\times6}{7\times6} + \frac{-13\times2}{21\times2} + \frac{5\times7}{6\times 7} \\
&=& \frac{24}{42} + \frac{-26}{42} + \frac{35}{42}\\
&=& \frac{24-26+35}{42}\\
&=& \frac{33}{42}\\
&=& \frac{3\times 11}{3\times 14}\\
&=& \frac{11}{14}
\end{eqnarray*}
\it \begin{eqnarray*}(-12)+(-11,57)+(+9,87) &=& -23,57 + 9,87 \\ &=& -13,7\end{eqnarray*}
\it \begin{eqnarray*}
(-4,51)+\left(-\frac45\right)+\left(+\frac7{20}\right)
&=& -4,51 - \frac{4\times2}{5\times2} + \frac{7\times5}{20\times5}\\
&=& -4,51 - \frac{8}{10} + \frac{35}{100}\\
&=& -4,51 - 0,8 + 0,35 \\
&=& -5,31 + 0,35 \\
&=& -4,96
\end{eqnarray*} 
\it \begin{eqnarray*}(-10)+(-15)-(-30)+(+7)-(+12)&=&
-10 - 15 +30 + 7 - 12 \\
&=& 30 + 7 -10 -15 -12\\
&=& 37 - 37 \\
&=&0\end{eqnarray*}
\it \begin{eqnarray*}
(-7)-(-9)-(+17)+(-41)+(+1)&=& -7 + 9 - 17 -41 + 1\\
&=& 9+1 -7-17-41\\
&=& 10 - 65 \\
&=&-55
\end{eqnarray*}
\it \begin{eqnarray*} (+49)-(-63)+(-3)-(+4)&=& 49+63-3-4\\
&=& 112 - 7 \\
&=& 105\end{eqnarray*}

\it \begin{eqnarray*} (-0,7)-(-0,9)+(-13,5)-(+4)&=& -0,7 +0,9 - 13,5 - 5\\
&=& 0,9 - 0,7-13,5-4 \\
&=& 0,9 - 18,2\\
&=& - 17,3\end{eqnarray*}
\it \begin{eqnarray*}\left(-\frac12\right)+\left(-\frac34\right) - \left(+\frac25\right)-(-1) &=& 
-\frac{1\times10}{2\times10}+\frac{-3\times5}{4\times5} + \frac{-2\times4}{5\times4} + 1 \\
&=& \frac{-10}{20} +\frac{-15}{20} + \frac{-8}{20} +\frac{20}{20} \\
&=& \frac{-10-15-8+20}{20}\\
&=& \frac{-13}{20}\end{eqnarray*}
\it \begin{eqnarray*}\left(-\frac5{11}\right)+\left(-\frac7{22}\right)
-\left(-\frac4{33}\right) &+& (-1) -(+0,5)\\
&=&
-\frac5{11}-\frac7{22}
+\frac4{33} -1 -0,5\\ 
&=&
 \frac{-5\times6}{11\times6} - \frac{7\times3}{22\times3} +\frac{4\times2}{33\times2} - \frac{66}{66} - \frac{0,5\times66}{66} \\
&=& \frac{-30}{66}-\frac{21}{66}+\frac8{66}-\frac{66}{66}-\frac{33}{66}\\
&=& \frac{-30-21+8-66-33}{66}\\
&=& \frac{-142}{66} \\ 
&=& \frac{-71}{33}
 \end{eqnarray*}
\it \begin{eqnarray*} 7-10+14-13-21+3 &=& 7 + 14 + 3 - 10 - 13 -21 \\
&=& 24 - 44 \\ &=& -20\end{eqnarray*}
\it \begin{eqnarray*}
0,75+4,7-0,7-0,5-1 &=& 5,45 - 2,2 \\ 
&=& 3,25
\end{eqnarray*}
\it \begin{eqnarray*}
\frac23-\frac34+\frac5{12}-\frac56 - 3 & = &
\frac{2\times4}{3\times4} - \frac{3\times3}{4\times3}
+\frac5{12} -\frac{5\times2}{6\times2} - \frac{3\times12}{12}\\
&=& \frac{8}{12}-\frac9{12} + \frac5{12} -\frac{10}{12} - \frac{36}{12}\\
&=& \frac{8-9+5-10-36}{12}\\
&=& \frac{-42}{12}\\
&=& \frac{-7\times6}{2\times6}\\
&=& -\frac72
\end{eqnarray*}
\it \begin{eqnarray*}[(5-9)+(3-5)] - [(7+3-5) - (7-10)]
&=& [-4 + -2] - [ 5 - (-3) ]\\
&=& -6 - [ 5 + 3] \\
&=& -6 - 8 \\
&=& -14\end{eqnarray*}
\it \begin{eqnarray*} [12-(14-5+0,75)] + [-15 + (3,25-2)]& =& [12 - 9,75] + [-15+ 1,25] \\
&=& 2,25 + (-13,75)\\
&=& -11,5\end{eqnarray*}
\it \begin{eqnarray*}\left[1 - \left(\frac25-\frac43\right)\right] &- &
\left[ \left(\frac45-1\right) - \left(\frac75+2\right)\right]\\&=& \left[ 1 - \left(\frac{2\times3}{5\times3}-\frac{4\times5}{3\times5}\right)\right] - \left[ \left(\frac45-\frac55\right) - \left(\frac75+\frac{2\times5}5\right)\right]\\
&=& \left[1 - \left(\frac6{15}-\frac{20}{15}\right)\right] - \left[ \frac{-1}5 - \left(\frac75+\frac{10}5\right)\right]\\
&=& \left[\frac{15}{15} - \frac{-14}{15}\right] - \left[\frac{-1}5-\frac{17}5\right]\\
&=& \left[\frac{15}{15}+\frac{14}{15}\right] - \left[\frac{-18}5\right]\\
&=& \frac{29}{15} +\frac{18}{15}\\
&=&\frac{47}{15}
\end{eqnarray*}

\end{itemize}
Supprimer les parenthèses et les crochets dans les expressions suivantes : 
\begin{itemize}
\it \begin{eqnarray*}(a-b+c) - (d- e - f) + (b-a) &=& 
a - b + c - d + e + f + b - a \\ 
&=& a - a - b + b + c - d + e + f \\
&=& c - d + e + f \end{eqnarray*}
\it \begin{eqnarray*}[(a-b)-(a-5)] + [b-7-(a-3)]&=&
(a-b)-(a-5) + b-7 -(a-3) \\ 
&=& a-b - a + 5 + b - 7 - a +3 \\
&=& a-a-a-b+b+5-7+3 \\
&=& - a +1 \end{eqnarray*}
ou : \begin{eqnarray*}[(a-b)-(a-5)] + [b-7-(a-3)]&=&
[a-b-a+5] + [b-7-a+3] \\
&=& [-b+5] + [-a+b-4] \\
&=& -b+5 -a+b-4 \\
&=& -a + b - b + 5 - 4 \\
&=& -a +1 
 \end{eqnarray*}
\it \begin{eqnarray*} [12-(a-b)+6] &-& [15+ (b-a-13)]\\&=&
12 - (a-b)+6 - 15 - (b-a-13)\\
&=& 12 - a + b + 6 - 15 - b + a + 13 \\
&=& -a + a +b - b +12 + 6 - 15 + 13\\
&=& 16
\end{eqnarray*}
ou :
\begin{eqnarray*} [12-(a-b)+6] &-& [15+ (b-a-13)]\\&=&
[12 - a + b + 6] - [15+b-a-13] \\
&=& [18-a+b]-[2+b-a] \\
&=& 18-a+b - 2 - b +a \\
&=& 18 -2  - a + a  + b - b \\
&=& 16	
\end{eqnarray*}
\end{itemize}

Effectuer les produits suivants : 
\begin{itemize}
\it \[(-3)(+7)(-5)(-4) = (- 21) \times (+20) = - 420\]
\it \[(+1,7)(-2,5)(-3)(-1,9) = (- 4,25)\times (+5,7) = -24,225\]
\it \[(-1)(+1)(-1)(+1)(-1)(-1) = (-1)^4 = +1\]
\it \[\left(-\frac35\right)(+3)\left(-\frac53\right) = 
\frac{-3 \times 3 \times -5}{5\times 3} = \frac{+3 \times (3\times5)}{3\times5} = 3\]
\it \[\left(+\frac12\right)\left(-\frac14\right)\left(-\frac48\right) = \frac{1\times -1\times -4}{2\times 4\times 8} = \frac{4}{2\times4\times8} = \frac1{2\times8} = \frac1{16}\]
\it \begin{eqnarray*}\left(+\frac12\right)\left(-\frac6{11}\right)\left(-\frac78\right)\left(-\frac{11}{15}\right)& =& \frac{1\times - 6 \times - 7\times -11} {2\times 11 \times 8 \times 15} \\&=& \frac{-2\times3\times7\times11}{2\times11\times2^3\times3\times5} \\&=&-\frac{7}{2^3\times5}\\&=&-\frac7{40}
 \end{eqnarray*}
\it \[\displaystyle(+15):\left(-\frac23\right)= 15 \times \frac{-3}2 = \frac{15\times -3}2 = -\frac{45}2\]
\it \[\displaystyle(-7,5):\left(-\frac45\right)= +7,5 \times \frac54 = \frac{7,5\times5}4= \frac{37,5}4 = \frac{37,5\times2}{4\times2} = \frac{75}8\]
\it \[\displaystyle(-1):\left(+\frac34\right) = -1\times \frac43 = -\frac43\]
\it \begin{eqnarray*}\displaystyle\left(-\frac45\right):(+1,4) &=& -\frac45: \frac{14}{10} \\&=& -\frac45\times\frac{10}{14} \\&=& - \frac{4\times 10}{5\times 14} \\&=& -\frac{2\times2\times2\times5}{5\times2\times7} \\&=& -\frac{2\times2}{7}\\&=& - \frac47\end{eqnarray*}
\it \[\left(+\frac34\right):\left(-\frac{15}8\right)= -\frac34\times8{15} = - \frac{3\times8}{4\times15} = - \frac{3\times2^3}{2^2\times3\times5} = - \frac{2}{5}.\]
\it \[\left(-\frac73\right):\left(+\frac53\right)= -\frac73\frac35=-\frac{7\times3}{3\times5} = -\frac75.\]
\it \[(+1):\left(+\frac23\right)= 1\times\frac32=\frac32\]
\it \[\left(-\frac{8}{25}\right):\left(-\frac{14}5\right)=+\frac8{25}\times{5}{14}=\frac{8\times5}{25\times14}=\frac{2^3\times5}{5^2\times2\times7}=\frac{2^2}{5\times7}=\frac4{35}\]
\it \[\left(+\frac13\right):\left(-\frac14\right) = 
- \frac13\frac41 = -\frac{1\times4}{3\times1}=-\frac43\] 
\it \begin{eqnarray*}\displaystyle\frac{-3+5-6}{2+7-1}\times\frac{4-3+1}{2-3-1}\times \frac{-5-7}{10}
&=& \frac{-4}{8} \times\frac2{-2}\times\frac{-12}{10}\\
&=& \frac{-4\times 2 \times -12}{8\times -2\times 10}\\
&=& \frac{2^2\times2\times2^2\times3}{-2^3\times2\times2\times5}\\
&=& -\frac{2^5\times3}{2^5\times5}\\
&=& - \frac35
\end{eqnarray*}
\it \begin{eqnarray*}\frac{\frac12-\frac23}{\frac34-1}\times \frac{\frac56+\frac13}{1-\frac45}\times \frac{-\frac78+1}{\frac32} &=& \frac{\frac{1\times3}{2\times3}
-\frac{2\times2}{3\times2}}{\frac34-\frac44}\times \frac{\frac56+\frac{1\times2}{3\times2}}{\frac55-\frac45}\times\frac{-\frac78+\frac88}{\frac32} \\
&=& \frac{\frac36-\frac46}{-\frac14}\times\frac{\frac56+\frac26}{-\frac15}\times\frac{\frac18}{\frac32}\\
&=&\frac{-\frac16}{-\frac14}\times\frac{\frac76}{-\frac15}\times\frac{\frac18}{\frac32}\\
&=& \left(-\frac16\right)\times(-4)\times\frac76\times(-5)\times\frac18\times\frac23\\
&=& - \frac{4\times7\times5\times2}{6\times6\times8\times3}\\
&=& - \frac{2^2\times7\times5\times2}{2\times3\times2\times3\times2^3\times3}\\
&=& - \frac{2^3\times5\times7}{2^5\times3^3}\\
&=& - \frac{5\times7}{2^2\times3^3} \\
&=& - \frac{35}{108}
\end{eqnarray*}
\it Calculer : \[ (-11)^3 =  1331; \phantom{miou}(+9)^2=81; \phantom{miou}(-10)^5=-100~000;\]\[ \phantom{miou}(+7)^4=2401; \phantom{miou}(-13)^4=28~561.\]
\it Calculer : \[ \left(-\frac12\right)^5 = \frac{(-1)^5}{2^5}=\frac{-1}{32}; \phantom{miou}\left(+\frac23\right)^2 = \frac{2^2}{3^2}=+\frac{4}9; \phantom{miou}\left(-\frac34\right)^3=(-1)^3\frac{3^3}{4^3}=-\frac{27}{64}\] \[ \phantom{miou}(+1,4)^2=+1,96; \phantom{miou}(-0,25)^2=+0,0625.\]
\it Calculer : \[ 5^3\times 5^2=5^5 =5^5 (=3~125) ; \phantom{miou}(-7)\times(-7)^2=(-7)^3 = -343;\]
\[ \phantom{miou}\left(\frac23\right)\times \left(\frac23\right)^3=\left(\frac23\right)^4=\frac{2^4}{3^4}=\frac{16}{81}; \]\[
\left[\left(-\frac14\right)^2\right]^2\left[(-2)^2\right]^3=
\left(\frac{+1}{4^2}\right)^{2}[4]^{3} = \frac1{4^4}4^3
=4^{-4}4^3=4^{-4+3}=4^{-1}=\frac14
\]
\end{itemize}
Réduire les expressions suivantes: 
\begin{itemize}
\it \begin{eqnarray*}
-\frac32 x + \frac54 x-3x^2+\frac{x}6&-&\frac52 x^2+5+4x^2\\&=& -3x^2-\frac52x^2+4x^2 -\frac32x+\frac54x+\frac16x+5\\
&=& \left(-3 - \frac52+4\right)x^2 +\left(-\frac32+\frac54+\frac16\right)x+5\\
&=& \left(1 - \frac52\right)x^2+\left(-\frac{3\times6}{2\times6}+\frac{5\times3}{4\times3}+\frac{1\times2}{6\times2}\right)x+5\\
&=& \left(\frac22-\frac52\right)x^2+\left(-\frac{18}{12}+\frac{15}{12}+\frac{2}{12}\right)x+5\\
&=& \frac{2-5}2x^2+\frac{-18+15+2}{12}x+5\\
&=&-\frac32x^2-\frac1{12}x+5
\end{eqnarray*}
\it \begin{eqnarray*}\frac32 x^2+xy+y^2-2xy&+&\frac{x^2}3-\frac32x^2\\
&=& \frac32x^2+\frac13x^2-\frac32x^2+y^2+xy-2xy\\
&=& \left(\frac32+\frac13-\frac32 \right)x^2+y^2+(1-2)xy\\
&=& \frac13x^2+y^2+(-1)xy\\
&=& \frac13x^2+y^2-xy
\end{eqnarray*}
\it \begin{eqnarray*}
4a^2-\frac23a-\frac35a^2&+&\frac13a-5a-\frac2{15}a^2\\
&=&4a^2-\frac35a^2-\frac2{15}a^2-\frac23a+\frac13a-5a\\
&=&\left(4-\frac35-\frac2{15}\right)a^2+\left(-\frac23+\frac13-5\right)a\\
&=& \left(\frac{4\times15}{15}-\frac{3\times3}{5\times3}-\frac2{15}\right)a^2 +\left(\frac{-2+1}3-\frac{5\times3}3\right)a\\
&=& \left(\frac{60}{15}-\frac9{15}-\frac2{15}\right)a^2
+\left(\frac{-1}3-\frac{15}3\right)a\\
&=& \frac{60-9-2}{15}a^2+\frac{-1-15}3a\\
&=& \frac{49}{15}a^2-\frac{16}3a
\end{eqnarray*}
\it \begin{eqnarray*}3x^2+\frac45 - \frac53x - 2x^2 &-&\frac35x^3 + 4 - 2 x^2 + 7x\\
&=& -\frac35x^3+3x^2-2x^2-2x^2-\frac53x+7x+\frac45+4\\
&=& -\frac35x^3+(3-2-2)x^2 +\left(-\frac53+7\right)x+\frac45+\frac{4\times5}5\\
&=& -\frac35x^3 +(-1)x^2 + \left(-\frac53+\frac{7\times3}3\right)x+\frac45+\frac{20}5\\
&=& -\frac35x^3-x^2+\left(-\frac53+\frac{21}3\right)x+\frac{4+20}5\\
&=& -\frac35x^3-x^2+\frac{-5+21}3x+\frac{24}5\\
&=& -\frac35x^3-x^2+\frac{16}3x+\frac{24}5
\end{eqnarray*}

\it \begin{eqnarray*} 4x^2 &-& \frac72 + \frac35x - \frac52x^2 + \frac43 x^3 - 5 + \frac32 x^3 + 7 - 2x
\\ &=&  \frac43 x^3+ \frac32 x^3 + 4x^2- \frac52x^2
+ \frac35x- 2x-\frac72-5+7\\
&=&  \left(\frac43 +\frac32\right) x^3 + \left(4- \frac52\right)x^2
+ \left(\frac35- 2\right)x-\frac72+2\\
&=&  \left(\frac{4\times2}{3\times2} +\frac{3\times3}{2\times3}\right) x^3 + \left(\frac{4\times2}2- \frac52\right)x^2
+ \left(\frac35- \frac{2\times5}5\right)x-\frac72+\frac{2\times2}2\\
&=&  \left(\frac{8}{6} +\frac{9}{6}\right) x^3 + \left(\frac{8}2- \frac52\right)x^2
+ \left(\frac35- \frac{10}5\right)x-\frac72+\frac{4}2\\
&=&  \left(\frac{8+9}{6}\right) x^3 + \left(\frac{8-5}2\right)x^2
+ \left(\frac{3-10}5\right)x+\frac{-7+4}2\\
&=&  \frac{17}{6} x^3 + \frac{3}2x^2
+ \frac{-7}5x+\frac{-3}2\\
&=&  \frac{17}{6} x^3 + \frac{3}2x^2
- \frac{7}5x-\frac{3}2\end{eqnarray*}
\it 
\begin{eqnarray*}\frac25 a^2b + 3a^3 &-& 4ab^2 + \frac52 a^2b + \frac72b^3 - b^3 + 2ab^2\\
&=& 3a^3+\frac25a^2b+\frac52a^2b-4ab^2+2ab^2+\frac72b^3-b^3\\
&=& 3a^3+\left(\frac25+\frac52\right)a^2b
+(-4+2)ab^2+\left(\frac72-1\right)b^3\\
&=& 3a^3 + \left(\frac{2\times2}{5\times2}+\frac{5\times5}{2\times5}\right)a^2b + (-2)ab^2+\left(\frac72-\frac22\right)b^3\\
&=& 3a^3+\left(\frac{4}{10}+\frac{25}{10}\right)a^2b
-2ab^2+\frac{7-2}2b^3\\
&=& 3a^3+\frac{29}{10}a^2b-2ab^2+\frac52b^3\end{eqnarray*}
 
\end{itemize}
Réduire : \begin{itemize}
\it \begin{eqnarray*}(3x-5)&+&[2x-5 - (3x-2y+4) - (4x-3y-9)]\\
&=& 3x-5 + [ 2x-5 - 3x+2y-4 - 4x+3y+9]\\
&=& 3x-5 + 2x-5-3x+2y-4-4x+3y+9\\
&=& 3x + 2x - 3x - 4 x + 2y + 3y -5 - 5 - 4 +9 \\
&=& (3+2-3-4)x+(2+3)y+ (-5-5-4+9)\\
&=& -2x +5y -5\end{eqnarray*}
\it \begin{eqnarray*}(2x-5y+7) &-& [(3x+2y-3)-(4x+4y-2)]-[2x-(3y+4)]\\
&=& 2x-5y+7-[3x+2y-3-4x-4y+2] - [2x-3y-4]\\
&=& 2x-5y+7 - 3x-2y+3+4x+4y-2-2x+3y+4\\
&=& 2x-3x+4x-2x-5y-2y+4y+3y+7+3-2+4\\
&=& (2-3+4-2)x+(-5-2+4+3)y+12\\
&=& 1x+0y+12\\
&=& x+12\end{eqnarray*}
\it \begin{eqnarray*}[(x-2y+5) &-& (3x+2y+7)]-[(2x+3)-(4y-2)]\\
&=& [x-2y+5-3x-2y-7]-[2x+3-4y+2]\\
&=& x-3x-2y-2y+5-7 - 2x-3+4y -2\\
&=& x - 3x -2x -2y - 2y +4y +5 - 7 - 3 - 2 \\
&=& (1-3-2)x+(-2-2+4)y -7\\
&=& -4x-7\end{eqnarray*}
\end{itemize}
Effectuer les produits suivants : 
\begin{itemize}
\it \begin{eqnarray*}(3a^2b^3)\left(\frac23ab^5\right)&=& 
3\times\frac23\times a^2a\times b^3b^5\\
&=& 2a^3b^8\end{eqnarray*}
\it \begin{eqnarray*}\left(\frac45a^3b^2c\right)\left(-\frac34abc^4\right)&=& \frac45\times\frac{-3}4\times a^3ab^2bcc^4\\
&=& \frac{-4\times3}{5\times4}a^{3+1}b^{2+1}c^{1+4}\\
&=& -\frac35a^4b^3c^5
\end{eqnarray*}
\it \begin{eqnarray*}\left(\frac47a^2xy^3\right)\left(-\frac52a^3y^4\right)
&=& \frac47\times\frac{-5}2a^2a^3xy^3y^4\\
&=& -\frac{2\times2\times5}{7\times2}a^{2+3}xy^{3+4}\\
&=& -\frac{10}7a^5xy^7
\end{eqnarray*}
\it \begin{eqnarray*}\left(-\frac34x^2y\right)\left(+\frac35a^3y^5\right)&=&-\frac34\times\frac35 a^3x^2yy^5\\
&=& -\frac{3\times3}{4\times5}a^3x^2y^{1+5}\\
&=& - \frac9{20}a^3x^2y^6
 \end{eqnarray*}
\it \begin{eqnarray*}\left(\frac94a^4x^2y^3\right)\left(-\frac43ax^2\right)&=& -\frac94\times\frac43a^4ax^2x^2y^3\\
&=&-\frac{4\times9}{4\times3}a^{4+1}x^{2+2}y^3\\
&=& - 3a^5x^4y^3
\end{eqnarray*}
\it \begin{eqnarray*}\left(\frac{14}3a^2b^3x\right)\left(-\frac67a^2b^5\right)&=& 
\frac{14}3\times\frac{-6}7\times a^2a^2b^3b^5x\\
&=& \frac{-14\times 6}{3\times7}\times a^{2+2}b^{3+5}x\\
&=& - \frac{2\times7\times2\times3}{3\times 7} \times a^4b^8x\\
&=& - 4a^4b^8x\end{eqnarray*}
\it \begin{eqnarray*}\left(-\frac72ax^2y\right)\left(-\frac8{15}b^3xy^2\right)\left(\frac5{21}abx^3\right)&=&
\left(\frac{-7}2\times\frac{-8}{15}\times \frac{5}{21}\right)aab^3bxx^3yy^2\\
&=& +\frac{7\times8\times5}{2\times15\times21}a^{1+1}b^{3+1}x^{1+3}y^{1+2}\\
&=& +\frac{7\times2^3\times5}{2\times3\times5\times3\times7}a^4b^4x^4y^3\\
&=& +\frac{2^2}{3^2}a^4b^4x^4y^3\\
&=& +\frac49a^4b^4x^4y^3
\end{eqnarray*}
\it \begin{eqnarray*}\left(-\frac23xy^2\right)^2(-4x^2y)&=& 
\left(-\frac23xy^2\right)\left(-\frac23xy^2\right)(-4x^2y)\\
&=& \left(\frac{-2}3\times\frac{-2}3\times\frac{-4}1\right)xxx^2y^2y^2y\\
&=&\frac{-2\times-2\times-4}{3\times3}x^{1+1+2}y^{2+2+1}\\
&=& -\frac{16}9x^4y^5
\end{eqnarray*}
\it \begin{eqnarray*}\left(\frac5{12}a^4b^2x\right)\left(-\frac27ax^2y^3\right)\left(-\frac{14}5b^2xy^4\right)&=& 
\left(\frac5{12}\times\frac{-2}7\times\frac{-14}5\right)a^4ab^2b^2xx^2xy^3y^4\\
&=& \frac{+5\times2\times2\times7}{2^2\times3\times7\times5}a^{4+1}b^{2+2}x^{1+2+1}y^{3+4}\\
&=& \frac13a^5b^4x^4y^7
\end{eqnarray*}
\it \begin{eqnarray*}\left(\frac35x^2y\right)^3\left(-\frac54xy\right)&=& 
\frac{3^3}{5^3}\left(x^2\right)^3y^3\times\left(\frac{-5}{2^2}xy\right)
\\
&=& \frac{3^3\times -5}{5^3\times2^2}x^{2\times3}y^3xy\\
&=& \frac{-3^3}{5^2\times2^2}x^{6+1}y^{3+1}\\
&=&- \frac{27}{100} x^7y^4\end{eqnarray*}
\end{itemize}
Effectuer les produits suivants.
\begin{itemize}
\it \begin{eqnarray*}(2x-7)(-3x+2)&=& 
2x\times(-3x)+2x\times2+(-7)\times(-3x)+(-7)\times2\\
&=& -6x^2+4x+21x-14\\
&=& -6x^2 + 25x-14\end{eqnarray*}
\it \begin{eqnarray*}&&(4x^5+7-2x^3)(x^3-2x)\\&=& 
4x^5 x^3+4x^5(-2x)
+7 x^3+7(-2x)+(-2x^3) x^3+(-2x^3)(-2x)\\
&=& 4x^8-8x^6+7x^3-14x-2x^6+4x^4\\
&=&4x^8-8x^6-2x^6+4x^4+7x^3-14x\\
&=&4x^8-10x^6+4x^4+7x^3-14x
\end{eqnarray*}
\it \begin{eqnarray*}(5x^3-2x)(3x-4x^2)
&=& 5x^33x+5x^3(-4x^2)+(-2x)3x+(-2x)(-4x^2)\\
&=& 15x^4-20x^5-6x^2+8x^3\\
&=& -20x^5 + 15x^4 +8x^3 -6x^2\\
\text{(Il est courant d'ordonner des}&&\text{ grandes aux petites puissances.)}\end{eqnarray*}
\it \begin{eqnarray*}&&(2x-7x^2+5x^3)(3x-5x^2+8)\\
&=& 2x3x+2x(-5x^2)+2x8+(-7x^2)3x+(-7x^2)(-5x^2)\\&& \ \ \ \ \ +(-7x^2)8+5x^33x+5x^3(-5x^2)+5x^38\\
&=& 6x^2-10x^3+16x
-21x^3+35x^4-56x^2+15x^3-25x^5+40x^3\\
&=& -25x^5 +35x^4-10x^3-21x^3+15x^3+40x^3+6x^2-56x^2+16x\\
&=& -25x^5+35x^4+24x^3-50x^2+16x
\end{eqnarray*}
\it \begin{eqnarray*}\left(-2x+\frac32\right)(4x+3)&=& 
-2x4x+(-2x)3+\frac32\times4x+\frac323\\
&=& -8x^2-6x+\frac{12}2x+\frac{3\times3}2\\
&=& -8x^2-6x+6x+\frac92\\
&=& -8x^2+\frac92\end{eqnarray*}
\it \begin{eqnarray*}\left(\frac83x-\frac32x^2+5\right)(4x^3-5x^2+7)&=& 
\frac83x\times4x^3 -\frac83x\times5x^2 +\frac83x\times7\\
&&-\frac32x^2\times4x^3 + \frac32x^2\times5x^2 -\frac32x^2\times7\\ 
&&+5\times4x^3-5\times5x^2 +5\times 7\\
&=& \frac{32}3x^4 -\frac{40}3x^3 +\frac{56}3x\\
&&-6x^5 +\frac{15}2x^4 - \frac{21}2x^2\\
&&+20x^3-25x^2+35\\
&=& -6x^5+\left(\frac{32}3+\frac{15}2\right)x^4+\left(
-\frac{40}3+20\right)x^3\\&&+\left(-\frac{21}2-25\right)x^2 +\frac{56}3x+35\\
&=&-6x^5 +\frac{109}6 x^4 + \frac{20}3x^3 - \frac{71}2x^2+\frac{56}3x+35
\end{eqnarray*}
\it \begin{eqnarray*}(7x^4-2x^3+4x^2)(3x^2-5)
&=& 7x^4\times 3x^2 -7x^4\times 5 \\
&& -2x^3\times 3x^2 +2x^3\times 5\\
&& +4x^2\times 3x^2 -4x^2\times 5\\
&=& 21x^6 - 35 x^4 - 6x^5 +10x^3 +12x^4 - 20x^2\\
&=& 21x^6 -6x^5 - 23 x^4 +10x^3 - 20x^2\end{eqnarray*}
\it \begin{eqnarray*}(2x^2-4x^3)(x^3-2x) &=& 2x^2\times x^3
-2x^2\times 2x -4x^3x^3 +4x^3\times 2x\\
&=& 2x^5 -4x^3 -4x^6 +8x^4\\
&=& -4x^6+2x^5+8x^4-4x^3
\end{eqnarray*}
\it \begin{eqnarray*}(2x^2-4+2x)(x^2+5-2x)
&=& 2x^2x^2 + 2x^2\times 5 -2x^2\times 2x\\
&& -4x^2 -4\times 5 +4\times 2x\\
&&+2xx^2+2x\times 5-2x\times2x\\
&=& 2x^4+10x^2-4x^3-4x^2-20\\&&+8x+2x^3+10x-4x^2\\
&=& 2x^4-4x^3+2x^3+10x^2\\&&-4x^2-4x^2+8x+10x-20\\
&=& 2x^4-2x^3+2x^2+18x-20\end{eqnarray*}
\it \begin{eqnarray*}\left(\frac54x^3-2x+\frac12\right)\left(\frac72x^3-\frac23x+x^2\right)&=& 
\frac54x^3\times\frac72x^3-\frac54x^3\times\frac23x+\frac54x^3x^2
\\
&&-2x\frac72x^3 +2x\times \frac23x-2xx^2\\
&&+\frac12\times\frac72x^3-\frac12\times\frac23x+\frac12x^2\\
&=& \frac{35}8x^6 -\frac56x^4 +\frac54x^5\\
&& -7x^4 +\frac43x^2-2x^3\\
&&+\frac74x^3-\frac13x+\frac12x^2\\
&=& \frac{35}8x^6 +\frac54x^5 -\left(\frac56+7\right)x^4 \\
&&+\left(\frac74-2\right)x^3 +\left(\frac43+\frac12\right)x^2-\frac13x\\
&=& \frac{35}8x^6 +\frac54x^5 - \frac{47}6x^4-\frac14x^3+\frac{11}6x^2-\frac13x
\end{eqnarray*}
\end{itemize}
 Calculer les produits suivants : 
\begin{itemize}
\it \begin{eqnarray*}(2x+3)(3x+2)(x-4)&=& (2x\times3x+2x\times 2 +3\times 3x+3\times2)(x-4)\\
&=& (6x^2+4x+9x+6)(x-4)\\
&=& (6x^2+13x+6)(x-4)\\
&=& 6x^2x -6x^2\times 4 +13xx -13x\times 4 +6x-6\times 4\\
&=& 6x^3 -24x^2 +13x^2-52x+6x-24\\
&=& 6x^3-11x^2-46x-24\end{eqnarray*}
\it \begin{eqnarray*}(5x-1)(2x+3)(7+4x) &=& 
(5x\times2x+5x\times3-1\times 2x-1\times 3)(7+4x)\\
&=& (10x^2+15x-2x-3)(7+4x)\\
&=& (10x^2+13x-3)(7+4x)\\
&=&(10x^2\times 7 + 10x^2\times4x
+13x\times 7 + 13x\times 4x
-3\times 7 - 3\times 4x)\\
&=& (70x^2 +40x^3+91x+52x^2-21-12x)\\
&=& (40x^3+122x^2+79x-21)
\end{eqnarray*}
\it \begin{eqnarray*}
(3x^2-1)(x+1)(x-1) 
&=& (3x^2-1)(x^2-1^2)\\
&=& (3x^2-1)(x^2-1)\\
&=& (3x^2\times x^2-3x^2\times 1 
- 1\times x^2 +1\times 1)\\
&=& (3x^4-3x^2-x^2+1)\\
&=& (3x^4-4x^2+1)
\end{eqnarray*}
\end{itemize}
Développer et réduire : 
\begin{itemize}
\it \begin{align*}5(3a^2-4b^3)&-[9(2a^2-b^3)-2(a^2-5b^3)]
\\
& = 5\times 3a^2 - 5\times 4b^3 - [9\times 2a^2 -9b^3 -2a^2+-2\times (-5b^3)]\\
& = 15 a^2 - 20b^3 - [18a^2-9b^3-2a^2+10b^3]\\
&= 15a^2-20b^3-[16a^2+b^3]\\
&= 15a^2-20b^3-16a^2 - b^3\\
&= -a^2 -21b^3
\end{align*}
\it \begin{align*}3a^2(2b-1)&-[2a^2(5b-3)-2b(3a^2+1)]
\\ & = 3a^2\times 2b +3a^2\times (-1) - [2a^2\times5b
-2a^2\times3-2b\times3a^2-2b\times1]
\\ &=6a^2b - 3a^2 - [10a^2b-6a^2-6a^2b-2b]\\
&=6a^2b-3a^2 -10a^2b+6a^2+6a^2b+2b\\
&=6a^2b-10a^2b+6a^2b-3a^2+6a^2+2b\\
&= 2a^2b+3a^2+2b
\end{align*}
\it \begin{align*}(2a+5b)(3a-2b)&-(2a-1)(3a+2b)-(a-2b)(5b-1)\\
=& 2a\times 3a -2a\times 2b + 5b\times 3a -5b\times 2b 
\\&- (2a\times 3a +2a\times2b-3a-2b)\\
&- (a\times5b-a\times1-2b\times5b+2b)\\
=&6a^2-4ab+15ab-10b^2-(6a^2+4ab-3a-2b)\\
&- (5ab-a-10b^2+2b)\\
=&6a^2+11ab-10b^2-6a^2-4ab+3a+2b\\
&-5ab+a+10b^2-2b\\
=& 6a^2-6a^2 +11ab-4ab-5ab-10b^2+10b^2\\
&+3a+a+2b-2b\\
=& 2ab+4a
\end{align*}
\it \begin{align*}(2x-3y)(5x-2y)&-(3x-2y)(2x+1)-(5x-y)(3y+1)\\
=&2x\times5x-2x\times2y-3y\times5x+3y\times2y\\
&- (3x\times2x+3x\times1-2y\times2x-2y\times1)\\
&- (5x\times3y+5x\times1-y\times3y-y\times1)\\
=&10x^2-4xy-15xy+6y^2\\
&-(6x^2+3x-4xy-2y)\\
&-(15xy+5x-3y^2-y)\\
=&10x^2-19xy+6y^2-6x^2-3x+4xy+2y \\
&- 15xy-5x+3y^2+y\\
=&10x^2-6x^2-19xy+4xy-15xy+6y^2+3y^2\\
&-3x-5x+2y+y\\
=& 4x^2-30xy+9y^2-8x+3y
\end{align*}

\it \begin{align*}(ax^2-b)(ax^2-2b)&+3b(ax^2-b)+b(b-1)
\\
&= ax^2ax^2-ax^2\times2b-bax^2+b\times2b + 3bax^2-3bb+bb-b\times1\\
&= a^2x^4-2abx^2-abx^2+2b^2+3abx^2-3b^2+b^2-b\\
&= a^2x^4-2abx^2-abx^2+3abx^2 +2b^2-3b^2+b^2-b\\
&= a^2x^4-0abx^2+0b^2-b\\
&=a^2x^4-b\end{align*}

\end{itemize}
Résoudre les équations suivantes : 
\begin{itemize}
\it $5(2x-3) - 4(5x-7)=19-2(x+11)$. 
\begin{sol}
\textbf{Résolution :}
On réduit : 
\begin{eqnarray*}
5\times2x-5\times3-4\times5x-4\times(-7) &=& 19 -2x-2\times11\\
10x-15-20x+28 &=& 19-2x-22\\
(10-20)x+28-15 &=& -2x+19-22\\
-10x+13 &=& -2x-3\end{eqnarray*}
Puis on applique les opérations habituelles deux deux côtés de l'égalité pour isoler $x$ : 
\begin{eqnarray*}
-10x+13+2x&=& -2x-3+2x\\
-8x+13&=& -3\\
-8x+13-13 &=& -3-13\\
-8x &=& -16\\
x &=& \frac{-16}{-8}
\end{eqnarray*}
ce qui montre que la seule solution possible est $x=+2$ après simplification de la fraction. 

\textbf{Vérification} : On vérifie que cette valeur est bien solution en calculant les deux termes de l'égalité initiale pour $x=+2$.
Le terme de gauche devient 
\begin{eqnarray*}
5(2x-3)-4(5x-7) &=& 5\left(2\times2-3\right) -4\left(5\times2-7\right)\\
&=& 5(4-3) - 4\left(10-7\right)\\
&=& 5\times 1- 4 \times 3\\
&=& 5-12 \\
&=& -7.\end{eqnarray*}
Le terme de droite devient 
\begin{eqnarray*}
19-2(x+11) &=& 19-2\left(2+11\right)\\
&=& 19- 2 \times13\\
&=& 19-26\\
&=& -7.\end{eqnarray*}
Ceci confirme que $x=2$ est solution de l'équation de départ.
\end{sol}
\it $4(x+3)-7x+17 = 8(5x-1)+166$. 
\begin{sol}
\textbf{Résolution :}
On réduit : 
\begin{eqnarray*}
4x+4\times 3-7x+17 &=& 8\times5x -8\times1 +166\\
4x+12-7x+17&=& 40x-8+166\\
(4-7)x+12+17 &=&40x+158 \\
-3x+29 &=&40x+158\end{eqnarray*}
Puis on applique les opérations habituelles deux deux côtés de l'égalité pour isoler $x$ : 
\begin{eqnarray*}
-3x+29-40x&=& 40x+158-40x\\
-43x+29&=& 158\\
-43x+29-29 &=& 158-29\\
-43x &=&129\\
x &=& \frac{129}{-43}
\end{eqnarray*}
ce qui montre que la seule solution possible est $x=-3$ après simplification de la fraction (car $129=43\times3$). 

\textbf{Vérification} : On vérifie que cette valeur est bien solution en calculant les deux termes de l'égalité initiale pour $x=-3$.
Le terme de gauche devient 
\begin{eqnarray*}
4(x+3)-7x+17 &=&  4(-3+3)-7\times(-3)+17\\
&=& 4\times 0 +21+17\\
&=&0+38\\
&=& 38.\end{eqnarray*}
Le terme de droite devient 
\begin{eqnarray*}
8(5x-1)+166&=& 8(5\times (-3)-1)+166\\
&=& 8(-15-1)+166\\
&=& 8\times (-16)+166\\
&=& -128+166\\
&=& 38.\end{eqnarray*}
Ceci confirme que $x=-3$ est solution de l'équation de départ.
\end{sol}
\it $17-14(x+1)= 13-4(x+1)-5(x-3)$. 
\begin{sol}
\textbf{Résolution :}
On réduit : 
\begin{eqnarray*}
17 - 14x - 14\times 1 &=& 13-4x-4\times 1 -5x+5\times (-3)\\
17-14x-14&=& 13-4x-4-5x+15\\
-14x+3&=& -9x+24 \end{eqnarray*}
Puis on applique les opérations habituelles deux deux côtés de l'égalité pour isoler $x$ : 
\begin{eqnarray*}
-14x+3+9x&=& -9x+24+9x\\
-5x+3&=& 24\\
-5x+3-3&=& 24-3\\
-5x&=& 21\\
x &=& -\frac{21}{5}
\end{eqnarray*}
ce qui montre que la seule solution possible est $x=-\frac{21}{5}$ (ou $x=-4,2$).

\textbf{Vérification} : On vérifie que cette valeur est bien solution en calculant les deux termes de l'égalité initiale pour $x=-4,2$.
Le terme de gauche devient 
\begin{eqnarray*}
17-14(x+1) &=&  17 - 14\left(-\frac{21}5+1\right)\\
&=& 17 - 14\times \frac{-16}5\\
&=& \frac{85}5 + \frac{224}5\\
&=& \frac{309}5.\end{eqnarray*}
Le terme de droite devient 
\begin{eqnarray*}
13-4(x+1)-5(x-3)
&=& 13-4\left(-\frac{21}5+1\right)-5\left(-\frac{21}5-3\right)\\
&=& 13-4\times \frac{-16}5 -5\times\frac{-36}5\\
&=& \frac{65}5 + \frac{64}5 +\frac{180}5\\
&=& \frac{309}5.
\end{eqnarray*}
Ceci confirme que $x=-\frac{21}5$ est solution de l'équation de départ.
\end{sol}

\it $5x+3,5+(3x-4) = 7x-3(x-0,5)$. 
\begin{sol}
\textbf{Résolution :} On commence par développer : 
\begin{eqnarray*} 5x + 3,5 + 3x- 4 &=& 7x - 3x +3\times0,5\\
8x -0,5 &=& 4x  + 1,5 \end{eqnarray*}
On isole alors $x$ avec les opérations habituelles : 
\begin{eqnarray*}
8x - 0,5+0,5 &=& 4x+1,5+0,5\\
8x &= &4x+2\\
8x-4x&=&4x+2-4x\\
4x &=& 2\\
x&=& \frac24
\end{eqnarray*}
ce qui montre que la seule solution possible est $x=\frac24=\frac12$. 

\textbf{Vérification} : On vérifie que cette valeur est bien solution en calculant les deux termes de l'égalité initiale pour $x=\frac12$.
Le premier membre devient \begin{eqnarray*}
5x+3,5+(3x-4) &=& 5\times0,5 + 3,5 + (3\times 0,5 -4) \\
&=& 2,5+3,5 + (1,5-4) \\
&=& 6 -2,5\\
&=& 3,5
\end{eqnarray*}
et le second membre devient 
\begin{eqnarray*}
7x-3(x-0,5) &=& 7\times\frac12 -3\times(\frac12-0,5) \\
&=& 7\times0,5 - 3\times 0 \\
&=& 3,5-0 \\
&=& 3,5.
\end{eqnarray*}
Ce qui donne le même résultat, donc $\frac12$ est bien solution. 
\end{sol}
\it $7(4x+3)-4(x-1)=15(x+0,75)+7$. 
\begin{sol}
\textbf{Résolution} : On développe : 
\begin{eqnarray*}
7\times 4x + 7\times 3 - 4x - 4\times (-1) &=& 15 x + 15\times 0,75 + 7 \\
28x + 21-4x +4 &=& 15x + 11,25+7\\
24x+25&=& 15x+18,25
\end{eqnarray*}
puis on isole $x$ : 
\begin{eqnarray*}
24x+25-15x&=& 15x+18,25-15x\\
9x+25 &=& 18,25\\
9x+25-25 &=& 18,25-25\\
9x &=& -6,75 \\
x &=& -6,75\div 9 \\
x &=& -0,75
\end{eqnarray*}
ce qui montre que la seule solution possible est  $-0,75$. 

\textbf{Vérification} : Vérifions que $x=-0,75$ est solution en calculant les deux termes. Le premier terme donne \begin{eqnarray*}
7(4\times(-0,75)+3)-4(-0,75-1) &=& 7(-3+3) -4\times(-1,75) \\
&=& 7\times 0 + 7\\
&=& 7,
\end{eqnarray*}et le second terme donne 
\begin{eqnarray*}
15(-0,75+0,75)+7 &=& 15\times 0 + 7 \\ &=& 7,
\end{eqnarray*}ce qui est identique, donc $x=-0,75$ est bien solution de l'équation.
\end{sol}
\it $17x+15(x-1)=-1-14(3x+1)$. 
\begin{sol}
On développe 
\begin{eqnarray*}
17x+ 15x-15 &=& -1 - 14\times 3x -14\\
32x-15 &=& -1-42x-14\\
32x-15 &=& -42x-15
\end{eqnarray*}
 et on isole $x$
\begin{eqnarray*}
32x-15+15 &=& -42x-15+15\\
32x &=& -42x\\
32x+42x&=& -42x+42x \\
74x&=& 0\\
x&=&\frac0{74}\\
x&=&0,
\end{eqnarray*}ce qui montre que la seule solution possible est $x=0$. 

\textbf{Vérification :}Pour $x=0$, le premier terme vaut 
\begin{eqnarray*}
17x+15(x-1)&=& 17\times 0 + 15(0-1) \\ 
&=& 0 + 15\times (-1)\\
&=& -15,
\end{eqnarray*} et le second terme 
\begin{eqnarray*}
-1-14(3x+1) &=& -1-14(3\times0+1)\\
&=& -1-14(0+1) \\
&=& -1 -14\times 1\\
&=& -1-14 \\
&=&-15,\end{eqnarray*}
ce qui est identique, donc $x=0$ est bien solution de l'équation. 
\end{sol}
\end{itemize}
Résoudre les équations suivantes (après développement, les termes en $x^2$ ou $x^3$ se simplifient) : 
\begin{itemize}
\it $(x-1)^2+(x+3)^2= 2(x-2)(x+1)+38$.
\begin{sol}
\textbf{Résolution} : On commence par développer : 
\begin{eqnarray*} x^2-2\times x \times 1 +1^2 + x^2 + 2\times 3x + 3^2 &=& 2(xx +x\times1-2 x -2\times 1) +38\\
x^2-2x+1+x^2+6x+9 &=& 2(x^2+x-2x-2)+38\\
x^2+x^2-2x+6x+1+9 &=& 2(x^2-x-2)+38\\
2x^2+4x+10 &=& 2x^2-2x-4 + 38\\
2x^2+4x+10 &=& 2x^2-2x+34
\end{eqnarray*}
On isole alors $x$ : 
\begin{eqnarray*}
2x^2+4x+10 - 2x^2= 2x^2-2x+34-2x^2\\
4x+10 &=& -2x+34\\
4x+10+2x &=& -2x+34+2x\\
6x+10 &=& 34\\
6x+10-10&=&34-10\\
6x&=& 24\\
x&=& \frac{24}6\\
x&=& 4,
\end{eqnarray*}ce qui montre que la seule solution possible est $x=4$.

\textbf{Vérification} : Calculons alors les deux membres de l'équation pour $x=4$. Le premier membre donne 
\begin{eqnarray*}(x-1)^2+(x+3)^2 &=& (4-1)^2+(4+3)^2 \\
&=& 3^2 + 7^2 \\
&=& 9 + 49 \\
&=& 58,\end{eqnarray*}
et le second membre donne \begin{eqnarray*}
2(x-2)(x+1)+38 &=& 2(4-2)(4+1)+38\\
&=& 2\times2\times5+38\\
&=& 2\times 10+38\\
&=& 20+38\\
&=& 58,\end{eqnarray*} ce qui est identique, donc $x=4$ est bien solution de l'équation.
\end{sol}
\it $5(x^2-2x-1)+2(3x-2)=5(x+1)^2$.
\begin{sol}
\textbf{Résolution} : On développe 
\begin{eqnarray*}
5x^2-5\times2x-5\times 1 + 2\times 3x-2\times 2 &=& 
5(x^2+2\times1\times x+1^2)\\
5x^2 - 10x - 5 + 6x - 4 &=& 5(x^2 + 2x+1)\\
5x^2 -4x - 9 &=& 5x^2+5\times 2x + 5\times 1\\
5x^2-4x-9 &=& 5x^2 + 10x + 5
\end{eqnarray*}
puis on isole $x$ 
\begin{eqnarray*}
5x^2-4x-9 -5x^2 &=& 5x^2+10x+5-5x^2\\
-4x-9  &=& 10x+5\\
-4x - 9 - 10x&=& 10x+5-10x\\
-14x-9 &=& 5 \\
-14x-9+9&=& 5+9 \\
-14x&=& 14\\
x&=& \frac{14}{-14}\\
x&=&-1,\end{eqnarray*} ce qui montre que la seule solution 
possible est $x=-1$.

\textbf{Vérification} : On calcule alors les deux membres pour $x=-1$. Le premier membre vaut 
\begin{eqnarray*} 5(x^2-2x-1)+2(3x-2)&=& 
5((-1)^2 -2(-1) -1) + 2(3\times(-1)-2) \\
&=& 5(1+2-1) + 2(-3-2) \\
&=&5\times 2 + 2\times -5 \\
&=& 10 -10 \\
&=&0\end{eqnarray*}
et le second membre \begin{eqnarray*}
5(x+1)^2 &=& 5(-1+1)^2 \\
&=& 5\times (0)^2\\
&=& 5\times 0 \\
&=&0\end{eqnarray*}
ce qui est identique, donc $x=-1$ est bien solution de l'équation.
\end{sol}
\it $(9x+1)(x-2)=(3x+4)(3x-5)$.
\begin{sol}
\textbf{Résolution} : On développe 
\begin{eqnarray*}
9x\times x +9x \times (-2) +1\times x +1\times (-2)
&=& 3x\times 3x + 3x\times(-5) + 4\times 3x +4\times(-5)\\
9x^2 -18x +x -2 &=& 9x^2 -15x +12x -20 \\
9x^2 - 17x -2 &=& 9x^2-3x -20,
\end{eqnarray*}
puis on isole $x$ 
\begin{eqnarray*}
9x^2-17x-2-9x^2 &=& 9x^2-3x-20 - 9x^2 \\
-17x-2 &=& -3x-20 \\
-17x-2 + 3x &=& -3x-20 +3x \\
-14x - 2 &=& -20 \\
-14x - 2 + 2 &=& -20 + 2 \\
-14x &=& -18 \\
x &=& \frac{-18}{-14} \\
x&=& \frac{9}7
\end{eqnarray*} ce qui montre que la seule solution possible est $x=\frac97$. 

\textbf{Vérification} : 
On calcule alors les deux membres pour $x=\frac97$ : le premier membre donne 
\begin{eqnarray*} 
(9x+1)(x-2) &=& \left(9\times\frac97+1\right)\left(\frac97-2\right)\\
&=& \left(\frac{81}7 + \frac77\right)\left(\frac97-\frac{14}7\right)\\
&=& \frac{88}7\times \frac{-5}7 \\
&=& \frac{-88\times 5}{7\times 7} \\
&=& - \frac{440}{49},
\end{eqnarray*} et le second membre 
\begin{eqnarray*}
(3x+4)(3x-5) &=& 
\left(3\times\frac97+4\right)\left(3\times\frac97-5\right)\\
&=& \left(\frac{27}7 + \frac{28}7\right)\left(\frac{27}7-\frac{35}7\right)\\
&=& \frac{55}7 \times \frac{-8}7 \\
&=& - \frac{55\times 8}{7\times 7} \\
&=& - \frac{440}{49},
\end{eqnarray*} ce qui est identique, donc $x=\frac97$ est bien solution de l'équation.\end{sol}
\end{itemize}
Résoudre les équations suivantes (on se ramènera au cas entier en multipliant par un nombre opportun).
\begin{itemize}
\it $\frac{x+5}4 - \frac{x-3}6 = \frac{x}3$. 
\begin{sol}
\textbf{Résolution :} On commence par éliminer les dénominateurs $4$, $6$ et $3$ en multipliant par un de leurs multiples communs, par exemple $12$ : 
\begin{eqnarray*} 12\left(\frac{x+5}4 - \frac{x-3}6 \right) &=& 12\times \frac{x}3\\
\frac{12(x+5)}4 - \frac{12(x-3)}6 &=& \frac{12x}3\\
\frac{12}4(x+5) - \frac{12}6(x-3) &=& \frac{12}3 x \\
3(x+5) - 2(x-3) &=& 4x \end{eqnarray*}
puis on développe 
\begin{eqnarray*}
3x+3\times 5 -2x -2\times (-3) &=& 4x\\
3x+15-2x+6 &=& 4x \\
x+21 &=& 4x
\end{eqnarray*}et on isole $x$
\begin{eqnarray*}
x+21-4x &= &4x-4x\\
-3x+21 &=& 0 \\
-3x+21-21 &=& 0-21\\
-3x &=& -21 \\
x&=& \frac{-21}{-3}\\
x&=& +7\end{eqnarray*}
ce qui montre que la seule solution possible est $x=+7$. 

\textbf{Vérification} : On calcule alors les deux membres pour $x=7$. Le premier membre vaut 
\begin{eqnarray*} \frac{x+5}4 - \frac{x-3}6
&=& \frac{7+5}4 -\frac{7-3}6 \\
&=& \frac{12}4 - \frac46 \\
&=& 3 - \frac{2}{3} \\
&=& \frac93-\frac{2}{3}\\
&=& \frac73\end{eqnarray*}
et le second membre vaut 
\begin{eqnarray*}\frac{x}3&=& \frac73,
\end{eqnarray*} ce qui est identique, donc $x=7$ est bien solution.
\end{sol}
\it $\frac{3x-7}2 + \frac{x+1}3 = -16$. 

\begin{sol}
\textbf{Résolution :} On commence par éliminer les dénominateurs $2$ et $3$ en multipliant par un de leurs multiples communs, par exemple $6$ : 
\begin{eqnarray*}
6\left(\frac{3x-7}2 + \frac{x+1}3\right)&=& 6\times (-16)\\
\frac{6(3x-7)}2 + \frac{6(x+1)}3 &=& -96\\
\frac62(3x-7) + \frac63(x+1) &=& -96\\
3(3x-7) + 2(x+1)&=& -96,
\end{eqnarray*} puis on développe
\begin{eqnarray*} 
3\times 3x-3\times 7 + 2\times x+2\times 1&=& -96\\
9x - 21 + 2x + 2 &=& -96 \\
11x -19 &=& -96
\end{eqnarray*}et on isole $x$
\begin{eqnarray*}11x-19+19 &=& -96+19 \\
11x &=& -77\\
x&=& -\frac{77}{11}\\
x&=& -7
\end{eqnarray*}
ce qui montre que la seule solution possible est $x=-7$. 

\textbf{Vérification} : On calcule alors les deux membres pour $x=-7$. Le premier membre vaut 
\begin{eqnarray*}\frac{3x-7}2 + \frac{x+1}3 
&=& {3\times (-7) - 7}2 + \frac{-7+1}3 \\
&=& {-21-7}2 +\frac{-6}3 \\
&=& \frac{-28}2 -2 \\
&=& -14-2\\
&=&-16
\end{eqnarray*}
et le second membre vaut toujours $-16$ ce qui est identique, donc $x=-7$ est bien solution.
\end{sol}
\it $x-\frac{x+1}3 = \frac{2x+1}5$. 
\begin{sol}
\textbf{Résolution :} On commence par éliminer les dénominateurs $3$ et $5$ en multipliant par un de leurs multiples communs, par exemple $15$ : 
\begin{eqnarray*}
15\left(x-\frac{x+1}3\right) &=& 15\times \frac{2x+1}5\\
15x -\frac{15(x+1)}3 &=& \frac{15(2x+1)}5 \\
15x - \frac{15}3(x+1)&=& \frac{15}5(2x+1) \\
15x - 5(x+1) &=& 3(2x+1)
\end{eqnarray*} puis on développe
\begin{eqnarray*} 
15x - 5x-5\times 1 &=& 3\times 2x+3\times 1\\
10x - 5 &=& 6x+3
\end{eqnarray*}et on isole $x$
\begin{eqnarray*}
10x-5-6x &=& 6x+3-6x \\
4x-5 &=& 3 \\
4x-5+5 &=& 3+5 \\
4x &=&8 \\
x &=& \frac84\\
x&=&2
\end{eqnarray*}
ce qui montre que la seule solution possible est $x=2$. 

\textbf{Vérification} : On calcule alors les deux membres pour $x=2$. Le premier membre vaut 
\begin{eqnarray*}x-\frac{x+1}3&=& 2-\frac{2+1}3\\
&=& 
2-\frac33 \\&=& 2-1 \\ &=&1
\end{eqnarray*}
et le second membre vaut 
\begin{eqnarray*}\frac{2x+1}5 &=& \frac{2\times 2+1}5 \\
&=& \frac{4+1}5\\
&=&\frac55\\
&=& 1
\end{eqnarray*} ce qui est identique, donc $x=2$ est bien solution.
\end{sol}
\it $\frac{7-3x}{12} + \frac34 = 2(x-2) + \frac{5(5-2x)}6$. 
\begin{sol}
\textbf{Résolution :} On commence par éliminer les dénominateurs $12$, $4$ et $6$ en multipliant par un de leurs multiples communs, par exemple $12$ : 
\begin{eqnarray*}
12\left(\frac{7-3x}{12} + \frac34\right) &=&12\left(2(x-2) + \frac{5(5-2x)}6\right)\\
 \frac{12(7-3x)}{12}+\frac{12\times3}4 &=& 12\times2(x-2)
+\frac{12\times5(5-2x)}6\\
7-3x + 9 &=& 24(x-2) + \frac{60}6(5-2x)\\
-3x+16 &=& 24(x-2) + 10(5-2x)
\end{eqnarray*} puis on développe
\begin{eqnarray*} -3x+16 &=& 24x-24\times 2 + 10\times 5-10\times 2x\\
-3x+16 &=& 24x-48+50 -20 x \\
-3x+16 &=& 4x+2 
\end{eqnarray*}et on isole $x$
\begin{eqnarray*}
-3x+16-4x &=& 4x+2-4x \\
-7x+16 &=& 2 \\
-7x+16-16 &=& 2-16\\
-7x &=& -14\\
x &=& \frac{-14}{-7}\\
x&=& 2
\end{eqnarray*}
ce qui montre que la seule solution possible est $x=2$. 

\textbf{Vérification} : On calcule alors les deux membres pour $x=2$. Le premier membre vaut 
\begin{eqnarray*}\frac{7-3x}{12} + \frac34&=& 
\frac{7-3\times 2}{12}+\frac{3\times3}{4\times3}\\
&=& \frac{7-6}{12}+\frac{9}{12}\\
&=& \frac{1}{12}+\frac{9}{12} \\
&=& \frac{10}{12}\\
&=& \frac56
\end{eqnarray*}
et le second membre vaut 
\begin{eqnarray*}2(x-2)+\frac{5(5-2x)}6 &=& 
2(2-2) + \frac{5(5-2\times2)}6\\
&=& 2\times 0 + \frac{5(5-4)}6 \\
&=& 0 + \frac{5\times 1}6 \\
&=& \frac56
\end{eqnarray*} ce qui est identique, donc $x=2$ est bien solution.
\end{sol}
\it $\frac{x}5 - \frac{3x-1}6 + \frac{3-x}4 = 0$.
\begin{sol}
\textbf{Résolution :} On commence par éliminer les dénominateurs $5$, $6$ et $4$ en multipliant par un de leurs multiples communs, par exemple $60$ : 
\begin{eqnarray*}
60\left(\frac{x}5 - \frac{3x-1}6 + \frac{3-x}4 \right)&=& 60\times 0 \\
\frac{60x}5 - \frac{60(3x-1)}6 + \frac{60(3-x)}4 &=& 0\\ 
\frac{60}5x - \frac{60}6(3x-1) + \frac{60}4(3-x)&=&0\\
12x - 10(3x-1) +15(3-x) &=& 0
\end{eqnarray*} puis on développe
\begin{eqnarray*} 12x - 10\times 3x-10\times (-1) +15\times 3-15x &=& 0\\
12x-30x +10 +45 - 15x &=& 0 \\
(12-30-15)x +55 &=& 0 \\
-33x +55 &= &0
\end{eqnarray*}et on isole $x$
\begin{eqnarray*}-33x+55-55&=& 0-55\\
-33x&=& -55 \\
x&=&\frac{-55}{-33}\\
x&=& \frac{5}3
\end{eqnarray*}
ce qui montre que la seule solution possible est $x=\frac53$. 

\textbf{Vérification} : On calcule alors les deux membres pour $x=\frac53$. Le premier membre vaut 
\begin{eqnarray*}
\frac{x}5 - \frac{3x-1}6 + \frac{3-x}4 &=& \frac{\frac53}5 
- \frac{3\times\frac53-1}6+\frac{3-\frac53}4\\
&=& \frac13 - \frac{5-1}6 + \frac{\frac93-\frac53}4\\
&=& \frac13 - \frac46 + \frac{\frac43}4 \\
&=& \frac13- \frac23 + \frac13 \\
&=& \frac{1-2+1}3\\
&=&\frac03 \\
&=&0
\end{eqnarray*}
et le second membre vaut toujours $0$ ce qui est identique, donc $x=\frac53$ est bien solution.
\end{sol}
\it $\frac{3(x+3)}4 + \frac12 = \frac{5x+9}3 - \frac{7x-9}4$. 
\begin{sol}
\textbf{Résolution :} On commence par éliminer les dénominateurs $4$, $2$ et $3$ en multipliant par un de leurs multiples communs, par exemple $12$ : 
\begin{eqnarray*}12\left(\frac{3(x+3)}4 + \frac12\right)
&=& 12\left(\frac{5x+9}3 - \frac{7x-9}4 \right)\\
\frac{12\times3(x+3)}4 + \frac{12}2 
&=& \frac{12(5x+9)}{3} - \frac{12(7x-9)}{4} \\
\frac{36}4 (x+3) + 6 &=& \frac{12}3(5x+9)- \frac{12}4(7x-9) \\
9(x+3) +6 &=& 4(5x+9) - 3(7x-9)
\end{eqnarray*} puis on développe
\begin{eqnarray*} 
9x+9\times 3 +6 &=& 4\times 5x+4\times 9 - 3\times 7x-3\times(-9) \\
9x + 27 + 6 &=& 20 x +36 -21x + 27\\
9x + 33 &=& -x + 63
\end{eqnarray*}et on isole $x$
\begin{eqnarray*}
9x + 33 +x&=& -x + 63+x\\
10x+33 &=& 63 \\
10x+33-33&=&63-33 \\
10x&=&30\\
x&=&\frac{30}{10}\\
x&=& 3
\end{eqnarray*}
ce qui montre que la seule solution possible est $x=3$. 

\textbf{Vérification} : On calcule alors les deux membres pour $x=3$. Le premier membre vaut 
\begin{eqnarray*}\frac{3(x+3)}4 + \frac12
&=& \frac{3(3+3)}4 + \frac12\\
&=& \frac{3\times6}4 + \frac12\\
&=& \frac{18}4 + \frac12\\
&=& \frac92+\frac12\\
&=& \frac{10}2\\
&=& 5
\end{eqnarray*}
et le second membre vaut 
\begin{eqnarray*}\frac{5x+9}3 - \frac{7x-9}4
&=& \frac{5\times 3+9}3 - \frac{7\times 3- 9}4 \\
&=& \frac{15+9}3 - \frac{21-9}4 \\
&=& \frac{24}3 - \frac{12}4 \\
&=& 8 - 3 \\
&=& 5
\end{eqnarray*} ce qui est identique, donc $x=3$ est bien solution.
\end{sol}
\it $\frac{2x-7}5 + \frac{x+11}2 = -4$. 
\begin{sol}
\textbf{Résolution :} On commence par éliminer les dénominateurs $5$ et $2$ en multipliant par un de leurs multiples communs, par exemple $10$ : 
\begin{eqnarray*}10\left(\frac{2x-7}5 + \frac{x+11}2\right) &=& -4\times 10 \\
\frac{10(2x-7)}5 + \frac{10(x+11)}2 &=& -40 \\
2(2x-7) + 5(x+11) &=& -40
\end{eqnarray*} puis on développe
\begin{eqnarray*} 
2\times 2x-2\times 7 + 5x+5\times 11 &=& -40\\
4x- 14 + 5x + 55 &=& -40\\
9x+41 &=& -40 
\end{eqnarray*}et on isole $x$
\begin{eqnarray*}9x+41-41 &=& -40-41 \\
9x&=&-81 \\
x&=& -\frac{81}9\\
x&=& -9
\end{eqnarray*}
ce qui montre que la seule solution possible est $x=-9$. 

\textbf{Vérification} : On calcule alors les deux membres pour $x=-9$. Le premier membre vaut 
\begin{eqnarray*}\frac{2x-7}5 + \frac{x+11}2
&=& \frac{2\times (-9)-7}5 + \frac{-9+11}2 \\
&=& \frac{-18-7}5 + \frac22 \\
&=& \frac{-25}5 + 1 \\
&=& -5+1\\
&=& -4
\end{eqnarray*}
et le second membre vaut toujours $-4$ ce qui est identique, donc $x=-9$ est bien solution.
\end{sol}
\it $\frac{2x-3}3 - \frac{x-3}6 = \frac{4x+3}4 - 17$. 
\begin{sol}
\textbf{Résolution :} On commence par éliminer les dénominateurs $3$, $6$ et $4$ en multipliant par un de leurs multiples communs, par exemple $12$ : 
\begin{eqnarray*}
12\left(\frac{2x-3}3 - \frac{x-3}6\right)
&=& 12\left(\frac{4x+3}4 - 17\right)\\
\frac{12(2x-3)}3 - \frac{12(x-3)}6 
&=& \frac{12(4x+3)}4 - 12\times 17 \\
4(2x-3) - 2(x-3) &=& 3(4x+3) - 204
\end{eqnarray*} puis on développe
\begin{eqnarray*} 
4\times 2x-4\times3 - 2x-2 \times(-3) &=& 3\times 4x+3\times 3 - 204\\
8x - 12 -2x + 6 &=& 12x + 9 - 204\\
6x - 6 &=& 12x - 195
\end{eqnarray*}et on isole $x$
\begin{eqnarray*}
6x - 6-12x &=& 12x - 195-12x\\
-6x - 6 &=& -195 \\
-6x-6+6 &=& -195+6\\
-6x &=& -189 \\
x &=& \frac{-189}{-6} \\
x &=& \frac{63}2
\end{eqnarray*}
ce qui montre que la seule solution possible est $x=\frac{63}2$. 

\textbf{Vérification} : On calcule alors les deux membres pour $x=\frac{63}2$. Le premier membre vaut 
\begin{eqnarray*}\frac{2x-3}3 - \frac{x-3}6
&=& \frac{2\times\frac{63}2-3}3 - \frac{\frac{63}2-3}6\\
&=& \frac{63-3}3 - \frac{\frac{63}2-\frac62}6\\
&=& \frac{60}3 - \frac{\frac{57}2}6\\
&=& 20 - \frac{57}{12} \\
&=& 20 - \frac{19\times4}{3\times4}\\
&=& 20 - \frac{19}4 \\
&=& \frac{80}4-\frac{19}4 \\
&=& \frac{61}4
\end{eqnarray*}
et le second membre vaut 
\begin{eqnarray*}
\frac{4x+3}4 - 17&=& 
\frac{4\times\frac{63}2+3}4 - 17 \\
&=& \frac{126+3}4-\frac{68}4 \\
&=& \frac{129-68}4\\
&=& \frac{61}4
\end{eqnarray*} ce qui est identique, donc $x=\frac{63}2$ est bien solution.
\end{sol}
\it $\frac{5x-3}4 - \frac{7x-5}9 = \frac{x+19}6$. 
\begin{sol}
\textbf{Résolution :} On commence par éliminer les dénominateurs $4$, $6$ et $9$ en multipliant par un de leurs multiples communs, par exemple $36$ : 
\begin{eqnarray*}
36\left(\frac{5x-3}4 - \frac{7x-5}9 \right)
&=& 36\times \frac{x+19}6\\
\frac{36(5x-3)}4 - \frac{36(7x-5)}9 &=& 
\frac{36(x+19)}6 \\
9(5x-3) - 4(7x-5)&=& 6(x+19) 
\end{eqnarray*} puis on développe
\begin{eqnarray*} 
9\times 5x-9\times 3 - 4\times 7x-4\times(-5)&=& 6x+6\times19\\
45x-27-28x+20 &=& 6x+114\\
17x -7 &=& 6x +114  
\end{eqnarray*}et on isole $x$
\begin{eqnarray*}
17x -7 -6x&=& 6x +114 -6x \\
11x-7 &=& 114\\
11x-7+7 &=& 114+7\\
11x &=& 121 \\
x&=& \frac{121}{11} \\
x&=& 11
\end{eqnarray*}
ce qui montre que la seule solution possible est $x=11$. 

\textbf{Vérification} : On calcule alors les deux membres pour $x=$. Le premier membre vaut 
\begin{eqnarray*}\frac{5x-3}4 - \frac{7x-5}9 
&=& \frac{5\times 11-3}4 - \frac{7\times 11-5}9 \\
&=& \frac{55-3}4 - \frac{77-5}9 \\
&=& \frac{52}4 - \frac{72}9 \\
&=& 13 - 8 \\
&=& 5
\end{eqnarray*}
et le second membre vaut 
\begin{eqnarray*}\frac{x+19}6&=&= \frac{11+19}6\\
&=& \frac{30}6\\
&=& 5
\end{eqnarray*} ce qui est identique, donc $x=$ est bien solution.
\end{sol}
\it $\frac{5x+1}8-\frac{x-1}3 = \frac{4(2x-3)}9$. 
\begin{sol}
\textbf{Résolution :} On commence par éliminer les dénominateurs $3$, $8$ et $9$ en multipliant par un de leurs multiples communs, par exemple $72$ : 
\begin{eqnarray*}
72\left(\frac{5x+1}8-\frac{x-1}3 \right)&=&
72\times \frac{4(2x-3)}9\\
\frac{72(5x+1)}8 - \frac{72(x-1)}3 &=& 
\frac{72\times4(2x-3)}9 \\
9(5x+1)-24(x-1) &=& 8\times4(2x-3)\\
9(5x+1)-24(x-1)&=& 32(2x-3)
\end{eqnarray*} puis on développe
\begin{eqnarray*} 
9\times 5x+9\times 1-24x-24\times(-1)&=& 32\times 2x-32\times 3\\
45x+9-24x+24&=& 64x-96\\
21x+33&=& 64x-96
\end{eqnarray*}et on isole $x$
\begin{eqnarray*}
21x+33-64x&=& 64x-96-64x\\
-43x+33 &=& -96 \\
-43x+33-33&=&-96-33\\
-43x &=& -129 \\
x &=& \frac{-129}{-43}\\
x&=& +3
\end{eqnarray*}
ce qui montre que la seule solution possible est $x=3$. 

\textbf{Vérification} : On calcule alors les deux membres pour $x=3$. Le premier membre vaut 
\begin{eqnarray*}\frac{5x+1}8-\frac{x-1}3 &=&
\frac{5\times3+1}8-\frac{3-1}3\\
&=&\frac{15+1}8-\frac23\\
&=& \frac{16}8-\frac23\\
&=& 2-\frac23\\
&=& \frac63-\frac23\\
&=& \frac43
\end{eqnarray*}
et le second membre vaut 
\begin{eqnarray*}\frac{4(2x-3)}9&=&
\frac{4(2\times3-3)}9 \\
&=&\frac{4(6-3)}9\\
&=& \frac{4\times3}{3\times3}\\
&=&\frac43
\end{eqnarray*} ce qui est identique, donc $x=3$ est bien solution.
\end{sol}
\it $\frac{2x-1}3 - \frac{5x+2}7 = x +13$. 
\begin{sol}
\textbf{Résolution :} On commence par éliminer les dénominateurs $3$ et $7$ en multipliant par un de leurs multiples communs, par exemple $21$ : 
\begin{eqnarray*}21\left(
\frac{2x-1}3 - \frac{5x+2}7 \right)&=&21(x +13)\\
\frac{21(2x-1)}3 - \frac{21(5x+2)}7 &=& 21(x+13)\\
7(2x-1) - 3(5x+2) &=& 21(x+13)
\end{eqnarray*} puis on développe
\begin{eqnarray*} 
7\times2x-7\times1 - 3\times 5x-3\times2 &=& 21x+21\times13\\
14x -7 - 15x - 6 &=& 21x + 273\\
-x - 13 &=& 21x + 273
\end{eqnarray*}et on isole $x$
\begin{eqnarray*}
-x - 13-21x &=& 21x + 273-21x\\
-22x - 13 &=& 273 \\
-22x-13+13 &=& 273+13\\
-22x&=& 286 \\
x &=& \frac{286}{-22}\\
x&=&-13
\end{eqnarray*}
ce qui montre que la seule solution possible est $x=-13$. 

\textbf{Vérification} : On calcule alors les deux membres pour $x=-13$. Le premier membre vaut 
\begin{eqnarray*}\frac{2x-1}3 - \frac{5x+2}7 &=&
\frac{2\times(-13)-1}3 - \frac{5\times(-13)+2}7\\
&=&
\frac{-26-1}3 - \frac{-65+2}7\\
&=&
\frac{-27}3 - \frac{-63}7\\
&=& -9 + 9 \\
&=&0
\end{eqnarray*}
et le second membre vaut 
\begin{eqnarray*}x +13&=&-13+13\\
&=&0
\end{eqnarray*} ce qui est identique, donc $x=-13$ est bien solution.
\end{sol}
\it $\frac{8x+2}5 - \frac{x-11}7 = \frac{5x-3}2 - \frac{3x-1}4$.
\begin{sol}
\textbf{Résolution :} On commence par éliminer les dénominateurs $5$, $7$, $2$ et $4$ en multipliant par un de leurs multiples communs, par exemple $140$ : 
\begin{eqnarray*}
140\left(\frac{8x+2}5 - \frac{x-11}7 \right)
&=& 140\left(\frac{5x-3}2 - \frac{3x-1}4\right)\\
\frac{140(8x+2)}5 - \frac{140(x-11)}7 
&=& \frac{140(5x-3)}2 - \frac{140(3x-1)}4 \\
28(8x+2) - 20(x-11) &=& 70(5x-3)-35(3x-1)
\end{eqnarray*} puis on développe
\begin{eqnarray*} 
28\times 8x+28\times2 - 20x-20\times(-11) &=& 70\times 5x-70\times 3\\
&&\text{           }-35\times 3x-35\times (-1)\\
224x + 56 - 20x + 220 &=& 350x - 210 - 105x + 35 \\
204x+276 &=& 245x- 175
\end{eqnarray*}et on isole $x$
\begin{eqnarray*}
204x+276 -245x&=& 245x- 175-245x\\
-41x+276 &=& -175 \\
-41x+276-276 &=& -175-276 \\
-41 x &=& -451 \\
x&=&\frac{-451}{-41} \\
x&=& 11
\end{eqnarray*}
ce qui montre que la seule solution possible est $x=11$. 

\textbf{Vérification} : On calcule alors les deux membres pour $x=$. Le premier membre vaut 
\begin{eqnarray*}\frac{8x+2}5 - \frac{x-11}7 &=&
\frac{8\times11+2}5 - \frac{11-11}7\\
&=&\frac{88+2}5 - \frac07\\
&=&\frac{90}5-0 \\
&=& 18
\end{eqnarray*}
et le second membre vaut 
\begin{eqnarray*}\frac{5x-3}2 - \frac{3x-1}4&=&
\frac{5\times11-3}2-\frac{3\times11-1}4\\
&=& \frac{55-3}2-\frac{33-1}4 \\
&=& \frac{52}2-\frac{32}4\\
&=& 26 - 8 \\
&=& 18
\end{eqnarray*} ce qui est identique, donc $x=11$ est bien solution.
\end{sol}
\end{itemize}
Résoudre les équations :
\begin{itemize}
\it $(x-1)(x+2)(x-3)=0$. 
\begin{sol}
\textbf{Résolution :}
On est en présence d'une équation produit.Toute solution est solution de l'équation d'une des trois équations $x-1=0$, $x+2=0$, et $x-3=0$. 

D'une part, $x-1=0$ revient à $x-1+1=0+1$, c'est-à-dire $x=1$. 

D'autre part, $x+2=0$ revient à $x=-2$. 

Enfin, $x-3=0$ revient à $x=3$.

On a donc trouvé que les seules solutions possibles sont $x=1$, $x=-2$ et $x=3$. 


\textbf{Vérification :} Pour $x=1$, $(x-1)(x+2)(x-3) = (1-1)(1+2)(1-3) = 0(1+2)(1-3) = 0$.

Pour $x=-2$, $(x-1)(x+2)(x-3) = (-2-1)(-2+2)(-2-3) = (-2-1)0(-2-3) = 0$.

Pour $x=3$, $(x-1)(x+2)(x-3) = (3-1)(3+2)(3-3) = (3-1)(3+2)0=0$.
\end{sol}

\it $(x-3)(x-4)(x-5)=0$. 
\begin{sol}
\textbf{Résolution :}
On est en présence d'une équation produit.Toute solution est solution de l'équation d'une des trois équations $x-3=0$, $x-4=0$, et $x-5=0$. 

D'une part, $x-3=0$ revient à $x=3$

D'autre part, $x-4=0$ revient à $x=4$. 

Enfin, $x-5=0$ revient à $x=5$.

On a donc trouvé que les seules solutions possibles sont $x=3$, $x=4$ et $x=5$. 


\textbf{Vérification :} Pour $x=3$, $(x-3)(x-4)(x-5) = (3-3)(3-4)(3-5)=0(3-4)(3-5)= 0$.

Pour $x=4$, $(x-3)(x-4)(x-5) = (4-3)(4-4)(4-5)=(4-3)0(4-5)= 0$.

Pour $x=5$, $(x-3)(x-4)(x-5) = (5-3)(5-4)(5-5)=(5-3)(5-4)0= 0$.

\end{sol}

\it $(2x+1)(x+1)(4x-3) = 0$.
\begin{sol}
\textbf{Résolution :}
On est en présence d'une équation produit.Toute solution est solution de l'équation d'une des trois équations $2x+1=0$, $x+1=0$, et $4x-3=0$. 

D'une part, la première se résout en 
\begin{eqnarray*}
2x+1-1 &=& 0-1 \\
2x &=& -1\\
x&=& -\frac12
\end{eqnarray*}

D'autre part, la seconde se résout en 
\begin{eqnarray*}
x+1 -1 &=&0-1 \\
x&=& -1
\end{eqnarray*}

Enfin,la dernière se résout en 
\begin{eqnarray*}
4x-3+3 &=& 0+3\\
4x &=& 3\\
x&=& \frac34
\end{eqnarray*}

On a donc trouvé que les seules solutions possibles sont $x=-\frac12$, $x=-1$ et $x=\frac34$. 


\textbf{Vérification :} En exercice. 
\end{sol}
\it $(2x+1)(x+4)(3x+1)=0$. 
\begin{sol}
\textbf{Résolution :}
On est en présence d'une équation produit.Toute solution est solution de l'équation d'une des trois équations $2x+1=0$, $x+4=0$, et $3x+1=0$. 

D'une part, la première se résout en 
\begin{eqnarray*}
2x+1-1 &=& 0-1 \\
2x &=& -1\\
x&=& -\frac12
\end{eqnarray*}

D'autre part, la seconde se résout en 
\begin{eqnarray*}
x+4 -4 &=&0-4 \\
x&=& -4
\end{eqnarray*}

Enfin,la dernière se résout en 
\begin{eqnarray*}
3x+1-1 &=& 0-1\\
3x &=& -1\\
x&=& -\frac13
\end{eqnarray*}

On a donc trouvé que les seules solutions possibles sont $x=-\frac12$, $x=-4$ et $x=-\frac13$. 


\textbf{Vérification :} En exercice. 
\end{sol}
\it $x(5x+1)(4x-3)(3x-4)=0$. 
\begin{sol}
\textbf{Résolution :}
On est en présence d'une équation produit.Toute solution est solution de l'équation d'une des quatre équations $x=0$, $5x+1=0$, $4x-3=0$, et $3x-4=0$. 

D'une part, la première a immédiatement comme seule solution $x=0$. 

D'autre part, la seconde se résout en 
\begin{eqnarray*}
5x+1 -1 &=&0-1 \\
5x&=& -1\\
x&=& -\frac15
\end{eqnarray*}
La troisième se résout en 
\begin{eqnarray*}
4x-3+3 &=&0+3 \\
4x&=& 3\\
x&=& \frac34
\end{eqnarray*}
Enfin,la dernière se résout en 
\begin{eqnarray*}
3x-4+4 &=& 0+4\\
3x &=& 4\\
x&=& \frac43
\end{eqnarray*}

On a donc trouvé que les seules solutions possibles sont $x=0$, $x=-\frac15$, $x=\frac34$ et $x=\frac43$. 


\textbf{Vérification :} En exercice. 
\end{sol}
\it $5x(3x-7)=0$.  
\begin{sol}
\textbf{Résolution :}
On est en présence d'une équation produit.Toute solution est solution de l'équation d'une des deux équations $5x=0$ et $3x-7=0$

D'une part, la première se résout en 
\begin{eqnarray*} 5x &=& 0 \\
\frac{5x}5 &=& \frac05\\
x&=&0\end{eqnarray*}

D'autre part, la seconde se résout en 
\begin{eqnarray*}
3x-7+7&=&0+7\\
3x &=& 7\\
x &=& \frac73
\end{eqnarray*}
On a donc trouvé que les seules solutions possibles sont $x=0$ et $x=\frac73$. 


\textbf{Vérification :} En exercice. 
\end{sol} 
\it $x^2-3x=0$. 
\begin{sol}
\textbf{Résolution :}On factorise par $x$ pour trouver \[x(x-3)=0,\]
qui est une équation produit.Toute solution est solution de l'équation $x=0$ ou de l'équation $x-3=0$. Cette dernière équation admet comme solution $x=3$. 

On a donc trouvé que les seules solutions possibles sont $x=0$ et $x=3$. 


\textbf{Vérification :} Pour $x=0$, $x^2-3x=0^2-3\times0=0-0=0$.

Pour $x=3$, $x^2-3x=3^2-3\times3=9-9=0$.
\end{sol}
\it $5x^2+8x=0$. 
\begin{sol}
\textbf{Résolution :}On factorise par $x$ pour trouver \[x(5x+8)=0,\]
qui est une équation produit. Toute solution est solution de l'équation $x=0$, ou de l'équation $5x+8=0$. Cette dernière équation revient à $5x=-8$, donc à $x=-\frac85$.

On a donc trouvé que les seules solutions possibles sont $x=0$ et $x=-\frac85$. 


\textbf{Vérification :} Pour $x=0$, $5x^2+8x=5\times 0^2+8\times0=5\times 0+0=0+0=0$.

Pour $x=-\frac85$, \begin{eqnarray*}5x^2+8x&=&5\left(-\frac85\right)^2 + 8\times \frac{-8}5\\
&=& 5 \frac{64}{25} - \frac{64}5 \\
&=& \frac{64\times 5}{5\times 5} - \frac{64}5\\
&=& \frac{64}5-\frac{64}5=0.\end{eqnarray*}

Les deux nombres $0$ et $-\frac85$ sont donc bien solution de l'équation.
\end{sol}
\it $4x^2-\frac{7x}3=0$. 
\begin{sol}
\textbf{Résolution :} On factorise par $x$ pour trouver \[x(4x-\frac73)=0,\]
qui est une équation produit. Toute solution est solution de l'équation $x=0$, ou de l'équation $4x-\frac73=0$. 
Cette dernière équation se résout ainsi : 
\begin{eqnarray*}
4x-\frac73 &=& 0 \\
4x-\frac73+\frac73 &=& 0+\frac73 \\
4x &=& \frac73 \\
x &=& \frac73\div 4\\
x&=& \frac7{12}\end{eqnarray*}

On a donc trouvé que les seules solutions possibles sont $x=0$ et $x=\frac7{12}$. 


\textbf{Vérification :} Pour $x=0$, $4x^2-\frac{7x}3
= 4\times 0^2 - \frac{7\times0}3 = 4\times 0 - \frac03 = 0-0=0$.

Pour $x=\frac7{12}$, \begin{eqnarray*}4x^2-\frac{7x}3&=&
4\left(\frac7{12}\right)^2 - \frac{7\times\frac{7}{12}}3 \\
&=& 4\times \frac{49}{144} - \frac{\frac{49}12}3 \\
&=& \frac{4\times 49}{4\times 36} - \frac{49}{12\times3}\\
&=& \frac{49}{36}-\frac{49}{36}\\
&=&0.
\end{eqnarray*}

Les deux nombres $0$ et $\frac7{12}$ sont donc bien solution de l'équation.
\end{sol}
\end{itemize}
	
    \section{Problèmes de réflexion plus avancés}\label{pbs}
\begin{enumerate}
\item Quel nombre faut-il ajouter aux deux termes de la fraction $\frac7{13}$ pour qu'elle devienne égale à $\frac23$ ? 
\begin{sol}
\textbf{Mise en équation :}
Notons $x$ le nombre cherché. L'énoncé devient
\[ \frac{7+x}{13+x}=\frac23.\]

\textbf{Résolution de l'équation :}
Par égalité des produits en croix, cela équivaut à 
\[ (7+x)3 = (13+x)2\]
Par distributivité, on trouve 
\[ 7\times3+x\times3 = 13\times2+x\times2,\]
c'est-à-dire \[ 21+3x=26+2x.\]
On retranche $2x$ aux deux termes : 
\[21+x = 26\]
Et on trouve, en retranchant 21 aux deux termes :  \[ x = 26-21 = 5.\]

\textbf{Vérification :} On peut alors vérifier que $\frac{7+5}{13+5} = \frac{12}{18} = \frac{2\times6}{3\times6} = \frac23.$
\end{sol}


 \item Quel nombre faut-il retrancher aux deux termes de la fraction $\frac{17}{23}$ pour obtenir une fraction égale à $\frac58$ ? 
 \begin{sol}
\textbf{Mise en équation :}
Notons $x$ le nombre cherché. L'énoncé devient
\[ \frac{17-x}{23-x}=\frac58.\]

\textbf{Résolution de l'équation :}
Par égalité des produits en croix, cela équivaut à 
\[ (17-x)8 = (23-x)5\]
Par distributivité, on trouve 
\[ 17\times8-x\times8 = 23\times5-x\times5,\]
c'est-à-dire \[ 136-8x=115-5x.\]
On ajoute $8x$ aux deux termes : 
\begin{eqnarray*}
136& =& 115-5x+8x\\&=& 115+3x
\end{eqnarray*} 
Et on trouve, en retranchant 115 aux deux termes : \begin{eqnarray*}
136-115 &=& 3x\\
21 &=& 3x\\
\frac{21}3&=&x\\
x&=&7
\end{eqnarray*}

\textbf{Vérification :} On peut alors vérifier que $\frac{17-7}{23-7} = \frac{10}{16} = \frac{2\times5}{2\times8} = \frac58.$
\end{sol}



 \item Quel nombre faut-il ajouter aux deux termes de la fraction $\frac58$ et retrancher aux deux termes de la fraction $\frac34$ pour obtenir deux fractions égales ? \begin{sol}
\textbf{Mise en équation :}
Notons $x$ le nombre cherché. L'énoncé devient
\[ \frac{5+x}{8+x}=\frac{3-x}{4-x}.\]

\textbf{Résolution de l'équation :}
Par égalité des produits en croix, cela équivaut à 
\[ (5+x)(4-x)=(8+x)(3-x).\]
Par distributivité, on trouve 
\[5\times4 -5x + 4x-x^2 = 8\times3-8x+3x-x^2 \]
c'est-à-dire \[ 20 -x - x^2 = 24 - 5x - x^2.\]
On ajoute $x^2$ aux deux termes pour trouver 
\[ 20 -x = 24 - 5x,\]
puis on ajoute $5x$ pour trouver \[ 20 -x+5x = 24,\]
c'est-à-dire \[20+4x=24\]
Et on trouve, en retranchant 20 aux deux termes,
\[4x= 24-20,\]
puis 
\[4x=4\]
et donc
\[x=\frac44 = 1.\]

\textbf{Vérification :} On peut alors vérifier que $\frac{5+1}{8+1} = \frac{6}{9} = \frac{2\times3}{3\times3} = \frac23$ et $\frac{3-1}{4-1} = \frac23$, donc on a bien \[\frac{5+1}{8+1}=\frac{3-1}{4-1}.\]
\end{sol}






 \item Une personne dispose de deux heures pour effectuer une promenade. Elle part en tramway à la vitesse moyenne de $12$ km à l'heure et revient à pied, à la vitesse moyenne de $4$ km à l'heure. À quelle distance du point de départ devra-t-elle quitter le tramway ? 
 \begin{sol}
 \textbf{Mise en équation :} Notons $x$ la distance recherchée en kilomètres. On parcourt donc une fois la distance $x$ à $v_1=12$ km/h, puis une seconde fois la même distance à $v_2=4$ km/h. \\
 À l'aller, le temps requis est $t_1=\frac{d}{v_1}=\frac{x}{12}$. \\
 Au retour, le temps requis est $t_2=\frac{d}{v_2}=\frac{x}4$. \\
 On cherche donc à avoir $t_1+t_2 = 2$ heures, c'est-à-dire : \[\frac{x}{12}+\frac{x}4= 2.\]
 \textbf{Résolution de l'équation :}
 On multiplie cette équation par $12$ pour se débarrasser des dénominateurs : 
 \begin{eqnarray*}
   12\left(\frac{x}{12}+\frac{x}4\right)&=&12\times2\\
   12\times\frac{x}{12}+12\times\frac{x}4 &=& 24\\
   \frac{12}{12}x+\frac{12}4x &=& 24\\
   x + 3x &=& 24\\
   4x&=&24
\end{eqnarray*}
On divise alors par $4$ pour trouver 
\[x=24\div4 = 6.\]On doit donc descendre après $6$ kilomètres. \\
 \textbf{Vérification :} Si on descend après $6$ kilomètres, le trajet aller prendra $\frac{6 \text{ km}}{12 \text{ km/h}} = 0,5$ h et le trajet retour prendra $\frac{6 \text{ km}}{4 \text{ km/h}} = 1,5$ h. Le temps total sera donc $0,5\text{ h} + 1,5\text{ h} = 2\text{ h}$, comme requis.
 
 \end{sol}
 \item Un épicier achète un certain nombre de bouteilles de vin pour $225$ euros. En vendant la bouteille $1,65$ euros, il gagnerait autant que ce qu'il perdrait s'il la vendait $1,35$ euros. Quel est le nombre des bouteilles de vins ? 
 \begin{sol}
 \textbf{Mise en équation :} On note $N$ le nombre de bouteilles de vin.  \\
 S'il vendait les bouteilles à $1,65$ euros, il toucherait $N\times1,65= 1,65N$ euros. \\
 S'il vendait les bouteilles à $1,35$ euros, il toucherait $N\times1,35=1,35N$ euros. \\
 L'énoncé nous dit donc que les différences $1,65N-225$ et $225-1,35N$ sont identiques, d'où \[1,65N-225=225-1,35N.\]
 \textbf{Résolution de l'équation :} On ajoute $1,35N$ aux deux termes : 
 \begin{eqnarray*}
 1,65N + 1,35N -225 &=& 225-1,35N+1,35N\\
 3N -225 &=& 225
 \end{eqnarray*}
 On ajoute alors $225$ aux deux termes : 
 \begin{eqnarray*}
 3N -225+225 &=& 225+225\\
 3N &=&450
 \end{eqnarray*}
 Et on divise par $3$ :  
 \begin{eqnarray*}
 \frac{3N}3 &=&\frac{450}3\\
 N &=& 150
 \end{eqnarray*}
 Il a donc acheté $150$ bouteilles de vin.

\textbf{Vérification :} S'il achète 150 bouteilles de vin, on peut vérifier que le montant total touché s'il vend à $1,65$ euros est de $150\times1,65 = 247,50$ euros. On réalise alors un bénéfice de $247,50-225 = 22,50$ euros.\\ 
S'il vendait à $1,35$ euros, il toucherait alors 
$150\times1,35 = 202,50$, ce qui représente une perte de $225-202,50=22,50$ euros. \\
Le bénéfice dans le premier cas est égal à la perte dans le second ce qui était la condition voulue.
 \end{sol}
 \item On considère un polygone régulier à $6$ côtés.
 \begin{enumerate}
 \item Réaliser une figure.
 \item Relier tous les sommets au centre $O$. 
 \item Montrer que tous les triangles ainsi formés ont leurs angles tous égaux à $60^o$. 
 \item En déduire que les angles d'un hexagone régulier mesurent $120$ degrés.
\end{enumerate}   

 \item On considère un polygone régulier à $5$ côtés.
 \begin{enumerate}
 \item Réaliser une figure.
 \item Relier tous les sommets au centre $O$. 
 \item Montrer que tous les triangles isocèles ainsi formés ont leurs angles à la base tous égaux à $54^o$. 
 \item En déduire que les angles d'un pentagone régulier mesurent $108^o$ degrés.
\end{enumerate}   
\item Utiliser les deux exercices précédents pour généraliser à un polygone régulier avec un nombre quelconque $N$ de côtés. On montrera ainsi que les angles d'une telle figure valent tous $\frac{N-2}N\times 180^o$.
\item On trace deux points $A$ et $B$ et une droite $(d)$ sur une feuille. Expliquer comment construire le cercle passant par $A$ et $B$ dont le centre appartient à $(d)$. Existe-t-il toujours ? Y en a-t-il toujours un seul ?
\item On place trois points $A$, $B$, $C$ sur une feuille. À quelle condition passe-t-il un cercle par ces trois points ? Comment le construire ? 
\item On considère un quadrilatère convexe $ABCD$ dans lequel la médiatrice de 
$[AB]$ est également la médiatrice de $[CD]$. Soit $O$ le point où cette médiatrice coupe la médiatrice de $[BC]$. \begin{enumerate}
\item Montrer que les quatre points $A$, $B$, $C$ et $D$ sont situés sur un même cercle de centre $O$. 
\item Démontrer que les médiatrices de $[AD]$, $[AC]$ et $[BD]$ passent également par $O$. 
\end{enumerate}
\item Tracer un parallélogramme $ABCD$, et les bissectrices de ses quatre angles. Elles se croisent intérieurement au parallélogramme en quatre points $E$, $F$, $G$ et $H$. Que dire du quadrilatère $EFGH$ formé ? Le démontrer. 
\item On trace un triangle $ABC$ et on place les milieux $D$, $E$, $F$ de ses trois côtés. Montrer que le triangle $DEF$ est une réduction du triangle $ABC$ dont on déterminera le coefficient. 
\item Dans le même contexte que l'exercice précédent, on relie chaque milieu au sommet opposé, et on note $G$ le point d'intersection des trois segments ainsi formés. Montrer que l'homothétie de centre $G$ et de rapport $-2$ envoie chaque sommet du triangle $DEF$ sur un sommet du triangle $ABC$. 
\item Dans un triangle $ABC$, on trace la hauteur issue de $A$ qui rencontre $(BC)$ en $H$. On suppose $AB^2 = BH \times BC$. Montrer que $ABC$ est rectangle en $A$. 
\item On considère un quadrilatère $ABCD$. Montrer que les milieux des côtés forment un  parallélogramme.
 \item Trouver trois nombres entiers consécutifs (consécutifs : qui se suivent) dont la somme est égale à $57$. 
 \begin{sol}
 \textbf{Mise en équation :} Si on note $x$ le premier des trois nombres, les trois nombres sont $x$, puis $x+1$ et enfin $x+2$. On a donc \[x + (x+1) + (x+2) = 57\]
 
\textbf{Résolution  de l'équation :} On réduit l'expression obtenue qui devient : \[3x+3 = 57\]
 Puis, en retranchant 3 : 
 \[3x = 54\]
 Et, en divisant par 3 : 
 \[x= 18\]
 On a donc trouvé que le premier nombre doit être $18$. Les deux autres sont donc $19$ et $20$.\\
 \textbf{Vérification :} On vérifie la solution : $18$, $19$ et $20$ sont bien trois entiers consécutifs, et leur somme est $18+19+20= 57$, ce qui convient.
 \end{sol}
 \item La somme de cinq nombres impairs consécutifs est 85. Quels sont ces nombres ? 
  \begin{sol}
 \textbf{Mise en équation :} Si on note $x$ le premier des trois nombres, les cinq nombres sont $x$, $x+2$, $x+4$, $x+6$ et $x+8$. On a donc \[x + (x+2) + (x+4) + (x+6) + (x+8) = 85\]
 
\textbf{Résolution  de l'équation :} On réduit l'expression obtenue qui devient \[5x+2+4+6+8 = 85\]
c'est-à-dire \[5x+20 = 85\]
 Puis, en retranchant 20 : 
 \[5x = 65\]
 Et, en divisant par 5 : 
 \[x= 13\]
 On a donc trouvé que le premier nombre doit être $13$. Les quatre autres sont donc $15$, $17$, $19$ et $21$.\\
 \textbf{Vérification :} On vérifie la solution : $13$, $15$, $17$, $19$ et $21$ sont bien cinq entiers impairs consécutifs, et leur somme est $13+15+17+19+21= 85$, ce qui convient.
 \end{sol}
 \item Un père a 29 ans ; son fils a 5 ans ; dans combien d'années l'âge du père sera-t-il le triple de l'âge du fils ? 
 \begin{sol}
 \textbf{Mise en équation :} Notons $N$ le nombre d'années recherché. Après ce temps, le père a $29+N$ années, et le fils $5+N$ années. La condition demandée s'écrit donc : 
 \[29+N = 3(5+N)\]
 
\textbf{Résolution  de l'équation :} On dévéloppe pour avoir \[ 29 + N = 3\times5+3N\]
puis \[29+N = 15+3N\]
On retranche N des deux côtés : 
\[29 +N - N = 15 + 3N -N\]
ce qui donne \[29 = 15+2N\]
On retranche alors $15$ des deux côtés : 
\[29-15 = 15+2N-15\]
ce qui donne \[14 = 2N\]
Puis, en divisant par 2 : \[7 = N\]
On doit donc attendre $N=7$ années.\\
 \textbf{Vérification :} On vérifie la solution : après $7$ années, le père a $29+7=36$ ans, et le fils a $5+7=12$ ans. Comme $36=3\times12$, le père a bien le triple de l'âge de son fils.
 \end{sol}
 \item Un père a 41 ans ; ses trois enfants sont âgés de 7, 9 et 13 ans. Dans combien d'années l'âge du père sera-t-il égal à la somme des âges de ses enfants ? 
 \begin{sol}
 \textbf{Mise en équation :} 
 Notons $N$ le nombre d'années recherché. Après ce temps, le père a $41+N$ années, et ses fils ont respectivement $7+N$, $9+N$ et $13+N$ années. On veut donc avoir \[41+N= (7+N)+(9+N)+(13+N)\]
 
\textbf{Résolution  de l'équation :} On réduit l'expression obtenue qui devient 
\[41+N = 7+9+13+N+N+N\]
donc \[41+N= 29+3N\]
On retranche alors $N$ des deux côtés pour obtenir 
\[41+N-N=29+3N-N\]
c'est-à-dire \[41=29+2N\]
puis on retranche 29 des deux côtés pour obtenir 
\[41-29=29+2N-29\]
c'est-à-dire 
\[12=2N\]
puis, en divisant par $2$, on trouve enfin 
\[6=N\]
Il faut donc attendre $6$ ans.\\
 \textbf{Vérification :} On vérifie la solution :
 après $6$ années, le père aura $41+6=47$ ans, et ses fils auront respectivement $7+6=13$ ans, $9+6 = 15$ ans, et $13+6=19$ ans. La somme des âges filiaux sera donc $13+15+19 =  47$ ans, ce qui est bien l'âge qu'aura le père. 
 \end{sol}
 \item Trouver un nombre dont la somme des quotients par 5, 7 et 9 soit égale à 429. 
 \begin{sol}
 \textbf{Mise en équation :} Notons $x$ ce nombre. Les quotients demandés sont $\frac{x}5$, $\frac{x}7$ et $\frac{x}9$. La condition demandée est donc 
 \[\frac{x}5+\frac{x}7+\frac{x}9=429\]
 
\textbf{Résolution  de l'équation :} On multiplie toute l'égalité précédente par $5\times7\times9=315$ pour éliminer les dénominateurs : 
\[315\left(\frac{x}5+\frac{x}7+\frac{x}9\right) = 429\times315\]
D'où : 
\[ \frac{315x}5 + \frac{315x}7+\frac{315x}9 = 429\times315\]
Comme $315=5\times63 = 7\times45=9\times35$, cela devient 
\[63x+45+35x = 429\times315\]
puis \[ 143x = 429\times315\]
et enfin \[x= \frac{429\times315}{143}\]
Cependant, $429=143\times3$, donc on peut simplifier :
\[x = \frac{143\times 3\times315}{143}=3\times315=945\]
Le nombre recherché est donc forcément $945$.\\
 \textbf{Vérification :} Il reste à vérifier que $945$ est bien solution du problème (on a montré que c'était la seule solution possible, mais pas forcément que c'était une solution). On calcule les quotients de l'énoncé : \begin{itemize}
 \item Le quotient par $5$ est $945\div5 = 189$. 
 \item Le quotient par $7$ est $945\div 7 = 135$. 
 \item Le quotient par $9$ est $945\div 9=105$. 
 \end{itemize}
 La somme de ces trois nombres est $189+135+105 = 429$, ce qui est la condition demandée. On a bien \[945\div5+945\div7+945\div9=429\]
 \end{sol}
 \item Développer les produits suivants : 
 \[ (x-1)(x+1)\]
 \[(x-1)(x^2+x+1)\]
 \[(x-1)(x^3+x^2+x+1)\]
 En déduire une formule générale. En déduire que 
 \[ x^n+ x^{n-1}+\ldots + x^2 + x + 1 = \frac{x^{n+1}-1}{x-1}\]
 \begin{sol}
 On développe : 
 \begin{eqnarray*}
 (x-1)(x+1) &=& xx +x\times 1 -1\times x -1\times 1 \\
 &=& x^2+x-x-1 \\
 &=& x^2-1
 \end{eqnarray*}
 Puis : 
 \begin{eqnarray*}
 (x-1)(x^2+x+1) &=& xx^2-1\times x^2 + xx +x\times 1 -1\times x -1\times 1 \\
 &=& x^3-x^2 + x^2+x-x-1 \\
 &=& x^3-1
 \end{eqnarray*}
 Et : 
 
 \begin{eqnarray*}
 (x-1)(x^3+x^2+x+1) &=& xx^3-1\times x^2+xx^2\\&&-1\times x^2 + xx +x\times 1 -1\times x -1\times 1 \\
 &=& x^4-x^3+x^3-x^2 + x^2+x-x-1 \\
 &=& x^4-1
 \end{eqnarray*}
 
En général, on en tire : 
 \begin{eqnarray*}
 (x-1)&\times &(x^n+x^{n-1}+\ldots+x^2+x+1) \\&=& xx^n-1\times x^{n-1}+xx^{n-1}-1\times x^{n-1}\\&&+\ldots+xx^2-1\times x^2 + xx +x\times 1 -1\times x -1\times 1 \\
 &=& x^{n+1}-x^n+x^n-x^{n-1}+\ldots+x^3-x^2+x^2-x+x-1 \\
 &=& x^{n+1}-1
 \end{eqnarray*}
 En divisant par $x-1$ les deux termes de cette égalité, il vient : 
 \[ x^n+x^{n-1}+\ldots + x^2+x+1 = \frac{x^{n+1}-1}{x-1}\]

 \end{sol}
 \item \  [Utiliser l'exercice précédent] Conte classique : un sage demande à un roi de le payer en mettant une pièce sur la première case d'un échiquier, puis deux pièces sur la seconde, puis quatre pièces, et ainsi de suite sur les soixante-quatre cases en doublant à chaque nouvelle case le montant. Déterminer la somme reçue. En utilisant l'approximation $2^{10} \sim 1000$, en donner un ordre de grandeur. 
 \begin{sol} 
 La première case continent une pièce, la case 2 contient $2$ pièces, la case 3 contient $2\times2 = 2^2$ pièces, la case 4 contient $2\times 2^2 = 2^3$ pièces. On en déduit successivement que la case numéro $k$ contient $2^{k-1}$ pièces. Le montant total $N$ est alors : 
 \[N=1+2+2^2+2^3+\ldots + 2^63.\]
 L'exercice précédent permet de simplifier
 \[ N= \frac{2^{64}-1}{2-1} = 2^{64}-1.\] 
 En ordre de grandeur, on a $2^{64} = 2^4\times 2^{60} =2^4\times(2^{10})^6\sim 16\times 1000^6 = 16.10^{18}\sim 2\times10^{19}$.
 Il y a donc environ $2.10^{19}$ pièces au total sur l'échiquier (c'est-à-dire vingt milliards de milliards). 
 
 \textbf{Remarque :} Un calcul exact donnerait $N= 18~446~744~073~709~551~615$. L'erreur commise par notre ordre de grandeur est donc de l'ordre de $13\%$.
 \end{sol}
 \item Un cercle délimite deux régions du plan : l'intérieur et l'extérieur. Deux cercles peuvent en déterminer quatre. S'en convaincre par un dessin. Que dire du nombre maximal de régions délimitées par trois cercles ? Que se passe-t-il pour quatre ? Trouver une formule générale.
 \item Si on coupe un disque par une corde, on délimite deux régions. Une deuxième corde délimite alors quatre régions. Généraliser. 
 \item Un nombre entier $N$ est dit \emph{pratique} si chaque entier entre $1$ et $N$ peut être écrit comme une somme de diviseurs de $N$ (sans répétition). Par exemple, $6$ est pratique car ses diviseurs sont $1$, $2$, $3$ et $6$ et que : 
 \[ 1 = 1\] 
 \[ 2 = 2 \]
 \[ 3 = 1 + 2 \]
 \[ 4 = 1+3 \text{ (mais 2+2 ne serait pas autorisé)}\]
 \[ 5 = 2 + 3 \]
 \[ 6 = 6\]
 Déterminer les nombres pratiques jusqu'à $30$. \\
 Montrer que toute puissance de $2$ est pratique (Penser à l'écriture en base $2$). 
\end{enumerate} 
 
\end{document}